% Auto-generated by scripts/06_emit_structure.py
% Source: scans 151–174
% Section id: chapter/5
% Phase 3 deliverable. Footnotes (Phase 4), citations (Phase 5),
% and Wade-Giles → Pinyin (Phase 6) are not yet applied here.

\chapter{The Oracle-Bone Inscriptions as Historical Sources}
\label{ch:5}

% source: scan 151, printed 134
% intro-echo-dropped: The Oracle-Bone
% intro-echo-dropped: Inscriptions as

\section[Other Sources: A Preliminary Note]{Other Sources: A Preliminary Note}
\label{sec:chapters-ch05:5.1}

% source: scan 151, printed 134
Oracle-bone inscriptions are not the only contemporary written sources available to the historian of \pinyinterm{shang} China. Putative \pinyinterm{shang} writings on silk, wood, or bamboo, have not been found, but we do possess nondivination records on bronze, bone, pottery, jade, and stone.\footnote[1]{For an introduction to nonoracular \pinyinterm{shang} records, see \pinyinterm{chen-mengjia-spaced} (\citeyear{Chen1956Yin}), pp. 44-46; \pinyinterm{su-yinghui} (\citeyear{Su1957Zhongguo}), p. 268; \pinyinterm{zheng-dekun} (\citeyear{Cheng1960ShangChina}), pp. 185-191. See too \ref{ch:1} (see ch. 1), nn. 21 and 22.} Certain bronze inscriptions, which all date from late in the historical period,\footnote[2]{For the evolution of inscriptions on bronze, see Kane (\citeyear{Kane1973Chronological}). Since most of the bronzes scientifically excavated at \pinyinterm{xiaotun} lack inscriptions (ibid., pp. 337-339), and since, as she argues, the great majority of \pinyinterm{shang} bronzes inscribed with an ancestor dedication come from period V (ibid., pp. 340-342), bronze inscriptions provide no information for the first part of the historical period. It is possible that some of the longer \pinyinterm{shang} bronze inscriptions were the product of provincial foundries (ibid., p. 361, n. 35; cf. Barnard [\citeyear{Barnard1958Recently}], p. 36, n. 21); if so, they stand in useful contrast to the oracle inscriptions which were documents made at the center of the kingdom.} have been useful for studying such matters as the late \pinyinterm{shang} calendar\footnote[3]{E.g., \pinyinterm{dong-short} (\citeyear{Tung1948bYinli}), pp. 198–207; \pinyinterm{shima-short} (\citeyear{Shima1966Bokuji}), pp. 12, 15-16; \pinyinterm{xu-jinxiong} (\citeyear{Hsu1971Wuchung}), p. 5a.} and aspects of \pinyinterm{shang} ancestor worship and lineage organization;\footnote[4]{E.g., \pinyinterm{zhang-guangzhi} (\citeyear{Chang1973Tan}), p. 116.} they present their own problems of decipherment, authenticity, and dating.\footnote[5]{\pinyinterm{shang} bronze inscriptions are, as Britton (\citeyear{Britton1940Fifty}), p. 5, noted, ``brief but plentiful.'' It has been estimated, for example, that of the 1,295 \pinyinterm{shang} and Western \pinyinterm{zhou-dynasty} bronzes that record \pinyinterm{tiangan} ancestral names, 1,102 are of \pinyinterm{shang} date (\pinyinterm{zhang-guangzhi} [\citeyear{Chang1973Tan}], p. 116). For an introduction to the provenance of the scientifically excavated bronzes, see Kane (\citeyear{Kane1975Reexamination}), p. 99, n. 14. On the general question of forged bronze inscriptions (though his concern is mainly with \pinyinterm{zhou-dynasty} models), see Barnard (\citeyear{Barnard1959Remarks}), (\citeyear{Barnard1968Incidence}), (\citeyear{Barnard1974Mao}); cf. \pinyinterm{zhang-guangzhi} (\citeyear{Chang1974Urbanism}); Matsumaru (1976, ed.). Akatsuka (\citeyear{Akatsuka1955Kohon}) provides transcriptions, translations, and commentary for 102 \pinyinterm{shang} bronze inscriptions, but the work should be used with caution. For an introduction to the clan insignia on the bronzes, see Hayashi Minao (\citeyear{Hayashi1968Inshu}). For a comprehensive approach relating the inscriptions to decorative motifs, vessel shapes, and the like, see \pinyinterm{zhang-guangzhi} (\citeyear{Chang1970Shang}), (\citeyear{Chang1973Tan}). For other introductory works, see \ref{ch:3} (see ch. 3), nn. 13-15 and 31.} The other records have to be used on a catch-as-catch-can basis; they are few in number and, on the whole, uninformative. This situation may change as new finds are made, but at present the tens of thousands of oracle-bone inscriptions, repetitive and formulaic though they frequently are, are the best written records we possess for understanding \pinyinterm{shang} culture. It goes without saying that the historian will also rely on other archaeological evidence, such as settlement sites, housing foundations, burial practices, and industrial remains. Studies in art history dealing with the vessel shape and decor of scientifically excavated bronzes, to cite but one example, promise to tell us much about the evolution and extent of \pinyinterm{shang} culture. Ethnographic studies of comparable cultures may also suggest new ways in which the \pinyinterm{shang} evidence may be interpreted, but such information should only be used when it can be well supported by the \pinyinterm{shang} evidence itself. Nothing in what follows should be taken to imply that the oracle-bone inscriptions are to be used to the exclusion of other evidence. A scholar who has mastered the oracle bones has not mastered all the written sources of \pinyinterm{shang}; and all these written sources must be used, when it is appropriate to do so, in conjunction with the unwritten archaeological evidence and the insights to be gained from comparative studies. Some \pinyinterm{zhou-dynasty} and \pinyinterm{han-dynasty} traditions may also be used to throw light on the \pinyinterm{shang} situation, but their historicity must be evaluated with great care in each case.

\section[The Dangers of Generalization]{The Dangers of Generalization}
\label{sec:chapters-ch05:5.2}

In using the oracle-bone inscriptions as historical sources, we must first beware
the ``archival fallacy''\footnote[9]{On the archival fallacy, see, in particular, Gelb (\citeyear{Gelb1967Approaches}), p. 4; cf. Fischer (\citeyear{Fischer1970Historians}), pp. 104-130.} which confuses the archive with the world. The mental contours
discerned in the divination inscriptions, which are religious and formulaic documents,
may not have been typical of \pinyinterm{shang} mentality in other areas of life. But there are, I
believe, good grounds for using the inscriptions, with appropriate caution,\footnote[7]{E.g., Kane (1974/75). Barnard and Satō (\citeyear{BarnardSato1975Metallurgical}), which catalogs all the reports of sites from which pre-Eastern \pinyinterm{han-dynasty} metallurgical remains were scientifically excavated in China between 1929 and 1966, is an invaluable reference for this kind of study.} as a guide to
at least certain aspects of the \pinyinterm{shang} world view.\footnote[8]{To cite but one example of the value of comparative studies, many insights into the role that divination played in \pinyinterm{shang} culture can be gained by reading accounts of divination in other societies; The works of Evans-Pritchard (\citeyear{EvansPritchard1937Witchcraft}), Middleton (\citeyear{Middleton1960Lugbara}), Beattie (\citeyear{Beattie1967Divination}), Park (\citeyear{Park1967Divination}) and Turner (\citeyear{Turner1975Revelation}) are of particular value. For an anthropological appreciation of \pinyinterm{shang} divination, see Monroe (\citeyear{Monroe1974Ritual}). See too the references cited in \ref{ch:3} (see ch. 3), n. 34.}\footnote[6]{Particularly interesting are the sixty-six \pinyinterm{shang} symbols, some of which are ancestral to later \pinyintext{jiaguwen} forms, found on pottery fragments at \pinyinterm{wucheng}, \pinyinterm{qingjiang}, \pinyintext{jiangxi}, 1973-74 (WW [1975.7], pp. 56-60; for the significance of the symbols, and their periodization, consult \pinyinterm{tang-lan} [\citeyear{TangLan1975Jiangxi}] and \pinyinterm{zhao-feng} [\citeyear{Zhao1976Qingjiang}]). These may be compared with the symbols inscribed on pot fragments from \pinyinterm{gaocheng}, \pinyinterm{taixicun}, \pinyintext{hebei}, excavated in 1973 (Ji Yun [\citeyear{Ji1974Kaocheng}]). More finds such as these, which provide important evidence of cultural diffusion over long distances, will throw clearer light than we now have on the origins of \pinyinterm{shang} writing. For a hypothesis about the Neolithic origins of the Chinese writing system, see Ho (\citeyear{Ho1975Cradle}), pp. 223-235; but see too my forthcoming review in Harvard Journal of Asiatic Studies 37.2 (1977).}
% source: scan 153, printed 136
I say this for two reasons. First, because the topics covered by the divinations are
numerous and varied. It is true that the bulk of the inscriptions is concerned with the
content, nature, and timing of religious sacrifices. But many other divinations cover wider
areas of royal concern: weather, agriculture, hunting, warfare, travel plans, selection of
personnel, sickness, dreams, childbearing, and so on (sec.~\ref{sec:chapters-ch02:2.3.1}% was: (sec. 2.3.1)
). Many of these are subjects
that we would regard as falling within the secular province of the meteorologist, farmer,
hunter, doctor, psychoanalyst, obstetrician, and so on, but the presence of these topics on
the bones indicates that for the \pinyinterm{shang} ruler they had their supernatural dimension. It
seems that virtually all political questions were religious questions too, and therefore
legitimate topics-in fact necessary topics--for divination. The secular and the sacred
were not distinct; the secular was an aspect of the sacred. The possibility that nonoracular,
nonreligious sources would tell a more secular tale cannot be excluded. But all the indications we have-archaeological, traditional, and ethnographic, as well as inscriptional—
present a consistent and convincing picture of a theocratic polity in which matters of
state were equally, and perhaps by definition, matters of religion.\footnote[10]{The issue will be discussed further in Studies. For the way divination has dominated existence in other societies and the way a religion may be characterized by divination, see, for example, Evans-Pritchard (\citeyear{EvansPritchard1937Witchcraft}), pp. 258-368 (cf. Winch [\citeyear{Winch1964Understanding}], p. 311); \textcite[pp.~162-163]{Lin1963Chi}; Hultkrantz (\citeyear{Hultkrantz1968La}), pp. 75, 86; Roux and Boratav (\citeyear{Roux1968Divination}), p. 282. These and other anthropological studies suggest that the \pinyinterm{shang} may also have practiced nonpyromantic forms of divination; oracles, once accepted, are sought in many forms.} %10

Second, I believe that we are justified in using the inscriptions as a reliable and central
source for \pinyinterm{shang} concerns because it is the king who, either directly or by proxy, divined.
Unless we are willing to believe that the \pinyinterm{shang} kings were hypocrites, it is clear that the
topics were important to them and were felt to be important to their ancestors and retainers.
The king was a theocrat who presumably ruled in part by virtue of his extraordinary
ability to communicate with the ultrahuman powers. The bone records, therefore, are not
simply a discarded priestly archive. They are religious-political records of decision making,
incantation, reassurance, and communication at the highest level of theocratic government.
If a topic does appear in the bones, we are justified in thinking that it was a topic of central
concern to the \pinyinterm{shang} kings and-since there are likely to have been few atheists in a
Bronze Age theocracy-to the \pinyinterm{shang} state as a whole. And if a topic is viewed in a certain
metaphysical way in the bones, it is plausible to assume that this cast of mind expresses
something that was central and fundamental to \pinyinterm{shang} culture.\footnote[11]{On the identity between divinatory and normal modes of thinking in societies in which divination is commonly accepted, cf. Vernant (\citeyear{Vernant1974Parole}), p. 10. The theological assumptions underlying \pinyinterm{shang} divination, together with divination's religious and political role, will be discussed in Studies. For a preliminary assessment, see Keightley (\citeyear{Keightley1972New}) and (\citeyear{Keightley1975Legitimation}).} %11

This is not to say that there were not other concerns, that there were not other modes
of conceptualization. The oracle inscriptions do not permit us to delineate the entire core
of \pinyinterm{shang} culture; but they do permit us to probe a significant segment of that core. For
the gathering of the plastrons and scapulas, the making of the cracks, and the carving of
% source: scan 154, printed 137
the inscriptions required much effort; that effort alone indicates how greatly pyromancy
must have been valued. \pinyinterm{shang} divination records, like the ritual bronzes, were the product
of intense mental effort, attention to cultic detail, and concentration of economic and
bureaucratic resources. It is not out of the question that if the \pinyinterm{shang} kings had had to
choose one set of documents to transmit to future historians, they would have chosen the
ones we now have.
Even when we turn to the subject of \pinyinterm{shang} religion, however, the topic most fully
documented in the inscriptions, we must still admit that the oracle bones are not fully
adequate sources. The picture of \pinyinterm{shang} religion which the inscriptions record is limited
and abstract. It is possible that some areas of \pinyinterm{shang} religion may have involved no divination at all, and of them we would have no record. The divinations we do have are but a
catalog or prompt book for what was presumably a far richer curriculum of cultic practice.
The music, the mimes, the chants, the visits of the ancestral powers, the cry of the victims,
the smell, smoke, flame, and blood-all the dramatic details may be imagined, but they
are unrecorded. We know the names of rituals, we know that there were dances; but we
have little idea what actions or feelings they involved. Inevitably, therefore, our view of
\pinyinterm{shang} religion is monochromatic and flat, drained of emotion. \pinyinterm{shang} ancestor worship
was ordered, hierarchical. bureaucratic-these traits are well recorded. The inscriptions
give us a true picture of a skeleton, but they do not allow us to generalize about the flesh
and blood that once animated it.\footnote[12]{On the problems involved in defining the essence of religion, and in systematizing varieties of religious experience, see Clavier (\citeyear{Clavier1968Resurgences}).} %12

A colleague once asked me if I knew the significance of the taotie ``monster mask''
found on so many \pinyinterm{shang} bronzes. I did not. ``If you don't understand the taotie,'' I
was told, ``you cannot understand the \pinyinterm{shang}.'' I still do not understand the taotie;\footnote[13]{For an introduction to the taotie motif and its possible significance, see Waterbury (\citeyear{Waterbury1942Early}), pp. 1-3, 129; \textcite[pp.~28, 39, 56, 59--60, 119, 121]{Linduff1972Tradition}; \pinyinterm{zhang-guangzhi} (\citeyear{Chang1973bFood}), pp. 516-518; Itō (\citeyear{Chang1975Ancient}), pp. 351-367.} %13
 it
is one of the numerous enigmas which the inscriptions have not solved. Let the taotie
serve as a salutary reminder of our ignorance and the need for caution.

\section[The Nature of the Sample]{The Nature of the Sample}
\label{sec:chapters-ch05:5.3}

All historical materials are incomplete, but some are more incomplete than
others. If it be accepted in the \pinyinterm{shang} case that the divinations record certain key concerns
of the elite and that used with the proper caution they may be used to generalize about
certain aspects
s of elite life, a further question remains. To what extent may the inscriptions
be used to generalize about themselves? In other words, to what extent, and how faithfully,
does the present corpus of inscriptions represent the total number of inscriptions originally
made?
% source: scan 155, printed 138
Typicality
Preliminary research suggests that we may now have approximately 5 to 10
percent of the original corpus of bones and shells inscribed by the \pinyinterm{shang} (see appendix
3). The smallness of this figure-which, though a crude estimate, is probably of the right
order makes it particularly important to determine the extent to which the content of
our sample may or may not be representative of the content of the original corpus. There
is no certain way to resolve this issue, and all answers must be provisional until every
square meter of the area once occupied by the \pinyinterm{shang} is excavated. My own conclusion,
which is highly tentative, is that the range of topics we have (sec.~\ref{sec:chapters-ch02:2.3.1}% was: (sec. 2.3.1)
) gives an accurate
indication of the true range of \pinyinterm{shang} divination but that the relative weight given any one
topic at any one time is less likely to be accurate.\footnote[14]{Such a conclusion applies to the whole corpus of approximately 47,000 fragments that have so far been published. (I obtain this figure by adding the estimated 5,000 fragments in Menzies and White to \pinyinposs{dong-short} 1953 estimate of about 42,000 fragments published; see appendix 3, n. 2). It cannot be applied to individual collections, many of which reflecting perhaps the storage pit or pits from which they derived, as well as the accidents of preservation-show an atypical distribution of materials by subject matter, period, etc. To cite but one example, \pinyinterm{renwen} has virtually no t'ien, ``hunt,'' inscriptions from period II, even though such divinations were common at that time (S289.3-290.1); similarly, it has remarkably few hunt inscriptions from period V (cf. S293.4-297.1).} %14

I draw this conclusion primarily from the patterns which may be discerned when the
inscriptions are classified by period. Some of these patterns, involving the prefatory formula
(sec.~\ref{sec:chapters-ch04:4.3.1.7}% was: (sec. 4.3.1.7)
), the frequency of prognostications and verifications (secs.~\ref{sec:chapters-ch04:4.3.1.8}, \ref{sec:chapters-ch04:4.3.1.9}% was: (secs. 4.3.1.8, 4.3.1.9)
),
the use of crack notations (sec.~\ref{sec:chapters-ch04:4.3.1.11}% was: (sec. 4.3.1.11)
), and the presence or absence of complementary
charge pairs (sec.~\ref{sec:chapters-ch04:4.3.1.10}% was: (sec. 4.3.1.10)
), are discussed in this book; others, such as the evolution of
religious conceptions, will require further treatment.\footnote[15]{For a preliminary assessment, see Keightley (\citeyear{Keightley1973aShang}), pp. 13-25, 37-38; (\citeyear{Keightley1975Legitimation}), pp. 49-52. The evolution of divination practice and religious conceptions will be considered more fully in Studies.} %15
 It is always possible that the excavation of new caches could disrupt the patterns so far established. The fact, however, that
the approximately 47,000 published fragments we now have, drawn from a variety of
pits and areas in the \pinyinterm{anyang} region, reveal a large number of consistent trends suggests
the more probable alternative-that these finds are indeed representative when considered
as a whole. Barring the discovery of storage pits used by a new and as yet unknown group
of diviners (pre-\pinyinterm{wu-ding}, nonroyal, or noncourt, for example), or of large stores of inscriptions dealing with topics that are poorly represented in the present corpus, at least for
certain periods,\footnote[16]{See, e.g., sec. 4.3.1.12; cf. \ref{ch:4} (see ch. 4), n. 138.} %16
 I doubt that new finds of inscriptions, though they may provide us with
information in greater depth, will radically affect our views of the \pinyinterm{shang}.\footnote[17]{Note that the preliminary report about the inscribed fragments excavated in 1973 (KK [\citeyear{KK1975Anyang}], pp. 27--46) contained few surprises.} %17
 Our divination
materials, in short, have the appearance of a random, representative scatter.
The only significant uncertainty in this regard concerns the virtual absence of inscriptions about the systematic, five-ritual cycle of the New School in period III+IV.\footnote[18]{Cf. \ref{ch:2} (see ch. 2), n. 27. For the five rituals, see \ref{ch:4} (see ch. 4), n. 102.} %18
 Whether
% source: scan 156, printed 139
the disappearance of the New School rituals of IIb is real or apparent is uncertain. A large
number of period III + IV finds came from the 大連坑, approximately 370 meters
northeast of the village center, and from sector F, approximately 90 meters south of the
village center, on the very edge of \pinyinterm{xiaotun};\footnote[19]{For TLK, see \pinyinterm{chen-mengjia-spaced} (\citeyear{Chen1956Yin}), p. 143; Itō (\citeyear{Ito1959Anyo}), pp. 349-351, 356; \pinyinterm{shi-zhangru} (\citeyear{Shih1959Hsiao}), p. 321. For sector F, see \pinyinterm{dong-short} (\citeyear{Tung1933Chia}), pp. 355-356; \pinyinterm{chen-mengjia-spaced} (\citeyear{Chen1956Yin}), pp. 140-143; KK (\citeyear{KK1975Anyang}), pp. 27-46.} %19
 the failure of the III + IV inscriptions we
now have to record the sacrifice cycle might be due to the fact that the storage pit (or pits)
in which such inscriptions may lie- period III+IV analog of YH127, perhaps—is
yet to be found. The fact that we already have far more inscriptions from period III+IV
than from periods II or V (sec.~\ref{sec:chapters-ch05:5.3.2}% was: (sec. 5.3.2)
), however, emphasizes the uncertainty of these
speculations.
-a
If, with that one possible, and major, exception, the random nature of our sample is
provisionally accepted, it follows that valid generalizations may be made about the scope
of \pinyinterm{shang} divination. Generalizations about quantitative matters, such as the percentage
of sacrificial wealth offered to particular ancestors, or the proportion of divinations concerning \pinyintext{di-high-god}, the ancestors, and the nature powers, or which enemy tribes were more troublesome than others, can be made with less certainty since the number of inscriptions involved
is frequently small and emphases could be significantly altered by the discovery of a dozen
new inscriptions. Even so, the fact that we are working with a large number of inscriptions
(see appendix \ref{app:3} (appendix 3)) suggests that the relative weight given any one particular topic is
not likely to be grossly inaccurate. One can only keep in mind that such generalizations
are not fully reliable, being based on what may be only 5 to 10 percent of all the inscriptions
made.

\subsection[Typicality]{Typicality}
\label{sec:chapters-ch05:5.3.1}

Preliminary research suggests that we may now have approximately 5 to 10 percent of the original corpus of bones and shells inscribed by the \pinyinterm{shang} (see appendix 3). The smallness of this figure—which, though a crude estimate, is probably of the right order—makes it particularly important to determine the extent to which the content of our sample may or may not be representative of the content of the original corpus. There is no certain way to resolve this issue, and all answers must be provisional until every square meter of the area once occupied by the \pinyinterm{shang} is excavated. My own conclusion, which is highly tentative, is that the range of topics we have (sec.~\ref{sec:chapters-ch02:2.3.1}% was: (sec. 2.3.1)
) gives an accurate indication of the true range of \pinyinterm{shang} divination but that the relative weight given any one topic at any one time is less likely to be accurate.\footnote[14]{Such a conclusion applies to the whole corpus of approximately 47,000 fragments that have so far been published. (I obtain this figure by adding the estimated 5,000 fragments in Menzies and White to \pinyinposs{dong-short} 1953 estimate of about 42,000 fragments published; see appendix 3, n. 2). It cannot be applied to individual collections, many of which reflecting perhaps the storage pit or pits from which they derived, as well as the accidents of preservation—show an atypical distribution of materials by subject matter, period, etc. To cite but one example, \pinyinterm{renwen} has virtually no t'ien, ``hunt,'' inscriptions from period II, even though such divinations were common at that time (S289.3-290.1); similarly, it has remarkably few hunt inscriptions from period V (cf. S293.4-297.1).} %14


I draw this conclusion primarily from the patterns which may be discerned when the inscriptions are classified by period. Some of these patterns, involving the prefatory formula (sec.~\ref{sec:chapters-ch04:4.3.1.7}% was: (sec. 4.3.1.7)
), the frequency of prognostications and verifications (secs.~\ref{sec:chapters-ch04:4.3.1.8}, \ref{sec:chapters-ch04:4.3.1.9}% was: (secs. 4.3.1.8, 4.3.1.9)
), the use of crack notations (sec.~\ref{sec:chapters-ch04:4.3.1.11}% was: (sec. 4.3.1.11)
), and the presence or absence of complementary charge pairs (sec.~\ref{sec:chapters-ch04:4.3.1.10}% was: (sec. 4.3.1.10)
), are discussed in this book; others, such as the evolution of religious conceptions, will require further treatment.\footnote[15]{For a preliminary assessment, see Keightley (\citeyear{Keightley1973aShang}), pp. 13-25, 37-38; (\citeyear{Keightley1975Legitimation}), pp. 49-52. The evolution of divination practice and religious conceptions will be considered more fully in Studies.} %15
 It is always possible that the excavation of new caches could disrupt the patterns so far established. The fact, however, that the approximately 47,000 published fragments we now have, drawn from a variety of pits and areas in the \pinyinterm{anyang} region, reveal a large number of consistent trends suggests the more probable alternative—that these finds are indeed representative when considered as a whole. Barring the discovery of storage pits used by a new and as yet unknown group of diviners (pre-\pinyinterm{wu-ding}, nonroyal, or noncourt, for example), or of large stores of inscriptions dealing with topics that are poorly represented in the present corpus, at least for certain periods,\footnote[16]{See, e.g., sec. 4.3.1.12; cf. \ref{ch:4} (see ch. 4), n. 138.} %16
 I doubt that new finds of inscriptions, though they may provide us with information in greater depth, will radically affect our views of the \pinyinterm{shang}.\footnote[17]{Note that the preliminary report about the inscribed fragments excavated in 1973 (KK [\citeyear{KK1975Anyang}], pp. 27-46) contained few surprises.} %17
 Our divination materials, in short, have the appearance of a random, representative scatter.

The only significant uncertainty in this regard concerns the virtual absence of inscriptions about the systematic, five-ritual cycle of the New School in period III+IV.\footnote[18]{Cf. \ref{ch:2} (see ch. 2), n. 27. For the five rituals, see \ref{ch:4} (see ch. 4), n. 102.} %18
Whether the disappearance of the New School rituals of IIb is real or apparent is uncertain.

\subsection[Chronological Distribution]{Chronological Distribution}
\label{sec:chapters-ch05:5.3.2}

We have no way of telling whether divination inscriptions were made in any
quantity prior to the reign of \pinyinterm{wu-ding}. The presence of two inscribed bone fragments at
\pinyinterm{zhengzhou},\footnote[20]{\pinyinterm{zhengzhou} \pinyinterm{erligang}, p. 38. The inscriptions---one thought to be a practice inscription---were on a bovid humerus and rib (\pinyinterm{chen-mengjia-spaced} [\citeyear{Chen1956Yin}], p. 6, plate 15; cf. \pinyinterm{li-xueqin-spaced} [\citeyear{LiHsuehchin1956Tan}], p. 17).} %20
 together with the fact that the \pinyinterm{wu-ding} inscriptions appear mature and
sophisticated in form, suggests that the royal generations immediately preceding \pinyinterm{wu-ding}
may have recorded their divinations on bone and shell. But few, if any, such records have
been discovered.\footnote[21]{See \ref{ch:4} (see ch. 4), n. 24.} %21

Within the historical era from \pinyinterm{wu-ding} to \pinyinterm{di-xin}, the inscriptions we have are not
distributed evenly by period. Preliminary research suggests that approximately 55 percent
of all fragments come from period I, 31 percent from III+IV, and less than 7 percent
% source: scan 157, printed 140
each from periods II and V (see table 31).\footnote[22]{These figures are crude in the extreme, particularly because the collections considered do not all treat the RFG inscriptions in the same manner. Nevertheless, the 35,208 fragments counted in table 31 represent about 74 percent of the fragments published; it is assumed, therefore, that the distribution is approximately correct.} Whether these variations are due to an accident
of discovery-so that large undiscovered caches of post-period I bones, especially those of
the New School of periods II and V, may still redress the balance-or whether it was primarily bones from these periods that the \pinyinterm{xiaotun} villagers ground up for medicine
and fertilizer or threw into dry wells is not certain.\footnote[23]{Cf. Lefeuvre (\citeyear{Lefeuvre1975Les}), p. 16.} %23
 Nor do we know whether these differences reflect the different durations of the kings' reigns and thus of the five periods,\footnote[24]{See the entries on fragments per year in table 31.} %24
 or
whether divination was indeed more frequent in period I than in all the other periods
combined. If we assume that the decreased scope and increased formalization of \pinyinterm{shang}
divination by period V (sec.~\ref{sec:chapters-ch04:4.3.1.12}% was: (sec. 4.3.1.12)
) was real and not apparent, the last alternative seems
more likely; it is plausible to assume that as divination became more routine, as it did
with the New School of period IIb, it may eventually have become more perfunctory and
less important, and would thus have been practiced less often.
Two conclusions follow from this uneven distribution. First, we are likely to know
more about \pinyinterm{shang} history in the reign of \pinyinterm{wu-ding} than in any other period. Second,
comparisons between activity in any two periods must always take into account the fact
that the size of the inscription sample varies by period; one inscription from period II,
for example, may carry the same relative weight as twelve from period I.

\section[The Problem of Authenticity]{The Problem of Authenticity}
\label{sec:chapters-ch05:5.4}

The carving of forged inscriptions, first on new bones and then on old, began as
soon as dealers recognized that collectors would pay by the graph.\footnote[25]{In the period when graphs were despised, the ancient bones sold for six cash for 0.59 kilograms; by the autumn of 1899, curio dealers were asking one china graph, and rich collectors were paying two and one-half liang of silver per graph; \pinyinterm{luo-zhenyu} paid nineteen silver dollars for the largest \pinyinterm{jinghua} fragment (Lefeuvre [\citeyear{Lefeuvre1975Les}], pp. 16, 18-19, 40, n. 1). The 韋-hsien dealers were making fakes from 1902 or 1903 on (ibid., p. 24, n. 1). See too Creel (\citeyear{Creel1938Studies}), p. 2; \pinyinterm{hu-houxuan} (\citeyear{Hu1955Yinhsu}), pp. 15-22; \pinyinterm{zhang-bingquan} (\citeyear{Zhang1967aTongming}), pp. 834-835.} %25
 The high prices offered
were a powerful incentive to embellish blank bone fragments-which constituted about
half the oracle bones found at \pinyinterm{xiaotun}\footnote[26]{For the proportion of blank to inscribed bones, see appendix 3, n. 5. Already provided with hollows, burn marks, and pu cracks, these blank bones were ideally suited to the forger's knife (cf. \pinyinterm{dong-short} [\citeyear{Tung1929bHsin}], p. 212).} %26
-or to supplement those already inscribed
with genuine inscriptions.\footnote[27]{The oracle bones reproduced as \pinyinterm{bali} 6-8, 10-12 (D) all bear fake and genuine inscriptions side-by-side (cf. n. 48), though \pinyinposs{rao-short} drawings do not reproduce the fake graphs. Fake and genuine inscriptions are also juxtaposed on \pinyinterm{kufang} 203, 983, 1598, 1604, 1619, and 1624 (all drawings) (\pinyinterm{jin-xiangheng} [\citeyear{Chin1962Shih}], pp. 189–191; \pinyinterm{li-yan-spaced} [\citeyear{LiYen1970PeiMei}], p. 283, plate 14) and 奇-chi W5 (see Britton's comment by the drawing). For forgeries inscribed on genuine bone, originally blank, see 奇-chi P76; P77; \pinyinterm{kufang} 2; 2175 (all drawings) (Britton [\citeyear{Britton1937aYin}], facing p. 1).} %27
 Fakes were produced in large numbers.\footnote[28]{``If one goes out to buy these bones he may encounter a hundred forgeries for every genuine piece. . . . I have seen collections of such bones, bought for many hundreds of dollars, not one of which was genuine'' (Creel [\citeyear{Creel1935Dragon}], p. 182). On the general prevalence of fakes at the time, with references to clandestine digs, nocturnal auctions, etc., see too Karlbeck (\citeyear{Karlbeck1957Treasure}), pp. 174--192. Large numbers of fakes were still on the market in Hongkong in 1946 to 1952 (\pinyinterm{li-yan-spaced} [\citeyear{LiYen1966Yenchai}], p. 3). See too n. 35.} %28
 In the view of Noel
% source: scan 158, printed 141
Barnard, ``With only a small supply of inscribed bones and plastrons available before the
large-scale diggings later commenced and as not a great deal was known about the script
or its contents, there existed,'' as Creel had noted, ``a veritable 'faker's paradise.''' And
he has made the further point that ``spurious inscriptions in genuine bone can be comparatively easily incised and made to appear quite authentic; when rubbings are taken
from them the deception may be hidden altogether.\footnote[29]{Barnard (\citeyear{Barnard1968Incidence}), p. 114, n. 15 (cf. Creel [\citeyear{Creel1938Studies}], pp. 2-3 Barnard (\citeyear{Barnard1960aReview}), p. 481. Barnard continues ([\citeyear{Barnard1960aReview}], p. 482) that he is not attempting ``to cast aspersions on the general corpus of unattested inscriptions; we may assume that a considerable proportion are quite acceptable.'' See too Barnard's comments quoted in n. 34.}''
 Great caution is in order. But, for the following reasons, the problem of forged inscriptions need not prove insurmountable.\footnote[30]{The prevalence of fakes, many of them crude, in the period shortly after the inscriptions caught the attention of scholars, encouraged conservatives like Zhang Taiyan (Binglin, d. 1936) who argued that all the inscriptions were fakes. (Uno Tetsuto [18751974], \pinyinterm{shima-kunio}'s teacher, also thought they were spurious [S608].) Zhang's arguments are notable mainly for their stubbornness; they fall into four main categories: \listitem{1} the early collectors were untrustworthy; \listitem{2} classical texts make no mention of inscriptions on shell (cf. \ref{ch:2} (see ch. 2), n. 155); \listitem{3} shells could not have lasted so long undecayed [the discovery of numerous skeletons from Yangshao sites invalidates this argument; the generally neutral pH of the soil around \pinyinterm{anyang} (Thorp [\citeyear{Thorp1936Geography}], p. 417, fig. 14) helps explain the preservation of bones, which are well preserved in basic soils with neutral or alkaline soil water (Goodyear [\citeyear{Goodyear1971Archaeological}], p. 147)]; \listitem{4} the inscriptions (which he at one point regarded as Northern Sung forgeries) were easy to fake (cited by \pinyinterm{dong-short} [\citeyear{Tung1965Chia}], pp. 56-60). Such skepticism still influenced Marcel Granet, who could write in 1929 (the date of the French edition of his book; that ``it would perhaps be an abuse of scepticism to refuse to trust in such labours as the learned M. \pinyintext{luo-zhenyu} has consecrated to the bones, which he has, no doubt, good reason to declare authentic'' (Granet [\citeyear{Granet1958Chinese}], p. 59). It is clear from the context that Granet felt the ``good reason'' to be partly commercial (cf. V. K. Ting [\citeyear{Ting1931ProfGranet}], pp. 288–289).} %30
 First, approximately 13,000 fragments published
in \pinyinterm{jiabian} and \pinyinterm{yibian} (some of which reappear in combined form in \pinyinterm{zhuihe} and \pinyinterm{pinbian})
were excavated by the Academia Sinica at \pinyinterm{xiaotun} between 1928 and\footnote[31]{For the number of fragments, see appendix 3, n. 2.} 1937.31 They
represent a corpus of authentic inscriptions against which others may be judged.\footnote[32]{In a few cases, I have been unable to determine the pit provenance of certain \pinyinterm{jiabian} inscriptions (see \ref{ch:4} (see ch. 4), n. 182). Until reports of their excavation are published, these inscriptions should not be regarded as fully attested.} %32
 Further,
as our knowledge of such factors as style, calligraphy, idioms, grammar, engraving and
erasing procedures, incised grooves, and inscription placement has developed, it has
become increasingly hard for the early fakes, with their inverted, misplaced graphs and
gibberish inscriptions, fabricated at a time when the knowledge of both forgers and scholars
was rudimentary, to escape detection.\footnote[33]{Yetts (\citeyear{Yetts1954Shang}), p. xxi, notes some of the anachronisms, including the use of small circles as punctuation marks, which caused Chalfant and his colleagues to assign a post-\pinyinterm{shang} date to the oracle bones. For one such spurious inscription, complete with the small circles, see \pinyinterm{kufang} 2175 (D) (cf. Britton [\citeyear{Britton1935Yin}], p. 21).} %33
 Authenticating criteria do exist.\footnote[34]{For an introduction to authenticating criteria, see White (\citeyear{White1945Bone}), pp. 75, 77; \pinyinterm{yan-yiping} (\citeyear{Yen1967Chiaku}), 1, pp. 2-6. \pinyinterm{zhou-hongxiang} (1967/68) and (\citeyear{Chou1970Jiagu}) has adumbrated some additional technical considerations that may, in the future, be employed in detecting fakes. For one reason why period V inscriptions are less likely to be faked, see Britton's remarks, quoted in \ref{ch:2} (see ch. 2), n. 102.} %34
 And these
% source: scan 159, printed 142
criteria have been used to identify a significant number of fakes in those collections which
were made in the early years of the twentieth century, when the collector's passion outran
his expertise.\footnote[35]{Western collectors generally received a paleographic education as they bought (cf. Yetts [\citeyear{Yetts1954Shang}], pp. xix-xx; \pinyinterm{yinxu-title}, preface, p. 5). The Reverend Frank Herring Chalfant wrote to L. C. Hopkins in 1910: ``My good wife and certain colleagues guy me about the business, and mockingly declare that there is a bone-factory in Weihsien city! I smile a grim smile, and am unmoved by this raillery.'' As Yetts (\citeyear{Yetts1954Shang}), p. xix, concludes, ``Their words in jest were only too true. Chalfant was being defrauded on a large scale. One less trustful might have taken warning from the city's notoriety for the faking of antiques.'' A total of 122 (7 percent) of the 1,687 fragments collected between 1904 and 1908 (?) (\pinyinterm{dong-short} [\citeyear{Tung1965Chia}], p. 61) and published in the Couling-Chalfant (\pinyinterm{kufang}) collection in 1935 are considered to be whole or partial fakes; they are listed, together with 4 more suspect pieces, by \pinyinterm{chen-mengjia-spaced} (\citeyear{Chen1956Yin}), pp. 652--653 (cf. Jung Keng [\citeyear{Jung1947Chiaku}], pp. 19-20). As \pinyinterm{chen-mengjia-spaced} notes, 114 of these 122 fragments were bovid bone, which the forger presumably found easier to carve than shell. James M. Menzies, who started to collect inscriptions in 1914 (\pinyinterm{yinxu-title}, preface, p. 1), probably bought more fakes than anybody else (he purchased the spuriously inscribed deer antler referred to in \ref{ch:1} (see ch. 1), n. 21), but eventually became a discriminating connoisseur (\pinyinterm{dong-short} [\citeyear{Tung1965Chia}], p. 61; \pinyinterm{zhang-bingquan} [\citeyear{Zhang1967aTongming}], p. 834; cf. Lefeuvre [\citeyear{Lefeuvre1975Les}], p. 56). The R. H. Lowie Museum of Anthropology at the University of California, Berkeley, contains three egregious fakes on bone (9.11970-2) that were evidently designed for the foreign market; only fifteen of the estimated one hundred fragments in the Britton collection in Columbia University's East Asiatic Library are thought to be genuine (\pinyinterm{li-yan-spaced} [\citeyear{LiYen1970PeiMei}], p. 320); the Musée Cernuschi contains at least four fakes (6093, 7232, 7630, 7691); there is a fake genealogical bone inscription (\pinyinterm{kufang} 1989 [D]) in the British Museum (Watson [\citeyear{Watson1963Handbook}], pp. 38-39; \pinyinterm{yan-yiping} [\citeyear{Yen1967Chiaku}], p. 5). Collections published by Chinese and Japanese scholars appear to contain few fakes; steeped in paleographic expertise, they were presumably able to exclude most fakes before publication. Among the fakes which have been identified in the early collections are Kikkō 2.28.12; 2.29.12; 2.29.16; \booktitle{Tieyun} 57.1; 84.1; 130.1; 254.1; \booktitle{Yechu} 26.1; 37.6; 37.8; \pinyinterm{yinxu-title} 758 (D) (\pinyinterm{dong-short} [\citeyear{Tung1953Chia}], p. 378; \pinyinterm{chen-mengjia-spaced} [\citeyear{Chen1956Yin}], p. 49; \pinyinterm{dong-short} [\citeyear{Tung1965Chia}], p. 61, also finds \booktitle{Tieyun} 256.1 a fake, but he appears to misread the last two graphs).} %35
 Noel
Barnard's ``principle of constancy of character
structures'' is rarely useful for assessing the authenticity of oracle-bone inscriptions. He describes the
principle thus: ``In each of the several thousands
of fully attested inscriptions I have studied it was
observed that repeated characters, or elements of
characters, in any one document, were always written
on the same structural principles--numbers of
strokes, position and intrinsic features of stroke
combinations, always accorded in each occurrence
of the same character written by the same writer''
([1958], p. 37; my italics). In authentic inscriptions, in short, ``structures do not vary although
shape differences are to be noted'' ([1959a], p. 25).
Barnard's qualification of the principle, however
\_``additions, omissions, or interchanges of elements in repeated characters . . . have nothing to
do with inconstancy but are regarded as a
perfectly normal mishap in writing'' (Barnard
[1973], p. 28; cf. p. 26, n. 8)-excludes from
consideration most of the epigraphic variations
found in the bone inscriptions upon which his
critics, frequently confusing shape with structure,
have seized with such delight (e.g., \pinyinterm{li-yan-spaced} [\citeyear{Yen1967Chiaku}];
\pinyinterm{zheng-dekun} [\citeyear{Cheng1971Inconstancy}]; for his spirited defense,
see Barnard [\citeyear{Barnard1965Chou}], pp. 408-424; [\citeyear{Barnard1973Chu}], pp.
23-28. Barnard himself has noted, in fact, that
``amongst unattested oracle bone texts about 20
cases of slight inconstancy were found in a total
of some 20,000 studied''; and in all but two of
these ``it was merely a matter of one or two strokes
omitted'' ([1959a], pp. 28-29). It is not this consideration, however, that renders the principle of
character constancy of little use-Barnard's 20,000
fragments, after all, may have been authentic.
The more serious problem is that the phrase ``in
any one document'' must have such restricted
meaning. It does not necessarily refer to the same
scapula or plastron, for the inscriptions found on
a single bone or shell were not necessarily written
at the same time (e.g., \pinyinterm{pinbian} 1; on this point,
see \ref{ch:3} (see ch. 3), n. 113); if not written at the same time,
inconstancy can easily be explained in terms of a
different engraver's hand; it has no bearing on
authenticity. Many of the counterexamples raised
by \pinyinterm{li-yan-spaced} (\citeyear{Yen1967Chiaku}), for example, can be dismissed
as not coming from the same document. The only
reasonable procedure is to assume that ``any one
document will usually mean a single preface
and charge, or, if appropriate, the two prefaces
and charges forming part of a charge pair (sec.~\ref{sec:chapters-ch02:2.5}% was: (sec. 2.5)
). But if only single charges are considered,
there will generally be no repetition of characters
by which character constancy may be tested
(Barnard [\citeyear{Barnard1973Chu}], p. 27, n. 10, makes the same
point). Paired charges do repeat many of the
characters, and here the principle does appear to
work, though certain curious exceptions may
weaken its usefulness. It is clear, for example, that
(a) and (b) are epigraphic variants of the
same graph and that the structure of the hand
element is significantly different. The fact that
both forms appear in the same attested document,
that is, in the same pair of inscriptions (\pinyinterm{pinbian}
134.5-6 [S98.1]; \pinyinterm{yibian} 8209 [S98.2]), violates
Barnard's principle. But the violation clearly accords with some \pinyinterm{shang} principle; that is, on these
two plastrons, form a appears in the right-hand
positive charge, form 6 in the left-hand negative
charge-though this distinction was not always
observed on other attested plastrons (on \pinyinterm{pinbian}
269.4-5, for example, form b was used for both
the positive and negative charges). A similar
violation of Barnard's principle, but with some
evident \pinyinterm{shang} rationale, may be found in the
case of (a) and (b) 3; form a is found on the
right side of the shell (\pinyinterm{pinbian} 269.1, 3, 5) and
form 6 on the left (\pinyinterm{pinbian} 269.4, 6) (cf. \pinyinterm{li-yan-spaced}
[1967], pp. 14-15). For other cases of what may
be intentional structural variation, see Serruys
(1974), p. 87. In short, the principle of constancy
of character structure has limited application for
the oracle bones, and when it does apply, it does
not do so invariably. It has, so far as I know, been
invoked but once-and then only in conjunction
with other criteria-in the identification of a
forgery (\pinyinterm{jin-xiangheng} [\citeyear{Chin1962aKufang}], p. 176).
% source: scan 160, printed 143
The possibility of skilled fakes must still be recognized.\footnote[36]{For the extensive career of one master forger, \pinyinterm{lan-baoguang}, who studied and copied from the works of Liu E and \pinyinterm{luo-zhenyu}, but without grasping the sense of the graphs or sentences he was forging, see \pinyinterm{dong-short} (\citeyear{Tung1965Chia}), p. 60.} %36
These are likely to be of four
types:
I. The forger copied an authentic inscription.\footnote[37]{\pinyinterm{lan-baoguang} (n. 36) advanced to the point where he would expertly copy whole fragments, a form, as \pinyinterm{dong-short} notes (1965, p. 60), of ``reprinting.''} %37
 If the copy is accurate, the bone,
and, more probably, its rubbing, may well escape detection. If so, no new, spurious information has contaminated the knowledge we gain from the inscriptions; judgments about
the frequency of a certain divination topic may be affected, but such judgments, as already
indicated (sec.~\ref{sec:chapters-ch05:5.3.1}% was: (sec. 5.3.1)
), are tentative at best.
2. The forger decided to create a new inscription. In that case, we may be dealing
with a hapax legomenon, which can readily be identified as such from the pages of \pinyinterm{shima-short}'s
\booktitle{Inkyo bokuji sōrui} (sec.~\ref{sec:chapters-ch03:3.3.2}% was: (sec. 3.3.2)
). Singularity is by no means a sign of forgery;\footnote[38]{For a classic case, see the disagreement between Britton and Menzies over fourteen fragments in the Fogg Art Museum at Harvard. As Britton concluded, ``The eccentric features simply raise the familiar puzzle, whether you have \listitem{1} a valuable rarity, or \listitem{2} a forgery'' (cited by \pinyinterm{zhou-hongxiang} [\citeyear{Chou1976OracleBone}], p. 21).} %38
 there are hundreds
of such unique inscriptions or graphs recorded in \booktitle{Inkyo bokuji sōrui}, many from attested collections. Generally speaking, no major conclusions should be drawn from any inscription
whose content exists in only the one exemplar and which is from a collection whose pieces
were excavated nonscientifically. The danger that the historian will be misled in such
cases is slight. First, because sound historical method would preclude giving any great
weight to a single inscription; if a topic were important to the \pinyinterm{shang} they presumably
divined about it frequently. Second, because it is frequently impossible to determine the
% source: scan 161, printed 144
meaning of many a hapax legomenon in the inscription record; only when a graph or formula
is repeated in a variety of contexts may we seize the meaning with confidence.\footnote[39]{\pinyinterm{dong-short} (\citeyear{Tung1951cChinese}), pp. 6-7, for example, neglected these considerations when he based his argument that the \pinyinterm{shang} calendar year was 365 days long on one inscription (\pinyinterm{yibian} 15 [S467.2]); the inscription is genuine but its meaning uncertain.} %39
3. The forger mass-produced a new inscription, thus avoiding the caveat of singularity
and lulling the historian, who finds safety in numbers, into a false sense of security. It is
wise to bear this possibility in mind.\footnote[40]{The Musée Guimet, for example, possesses more than a dozen fake inscribed scapula fragments which, from the similarity of their content and calligraphy, were all produced by the same modern atelier. The Metropolitan Museum of Art has at least one fake (67.43.16) that must have come from this same source. And I am informed by Christian Deydier that more inscriptions of this same type are in the British Museum.} %40
 Whenever we find a formula repeated in a series of
inscriptions which, from the history of the various collections in which the formula appears,
we have reason to think may have come into circulation about the same time,\footnote[41]{For the origin of the various collections, see \ref{ch:3} (see ch. 3), n. 4.} %41
 and whenever we find that this formula never appears in attested inscriptions from \pinyinterm{jiabian} or
\pinyinterm{yibian} or from their derivative collections, \pinyinterm{zhuihe} and \pinyinterm{pinbian}, we should at least be on
our guard.\footnote[42]{For example, since none are attested, the three inscriptions at S367.1 might all be fakes, with two copied from an original forgery (though, in this case, nothing about the rubbings gives ground for suspicion). The AC inscriptions at S356.2-3 provide an interesting case where three early, nonattested occurrences are vindicated by a later, attested one (\pinyinterm{zhuihe} 224, back).} %42
4. The forger decided to vary what he knew to be a genuine formula. This variation
should be treated with the cautions laid down in paragraph 2.
\subsection[Degrees of Reliability]{Degrees of Reliability}
\label{sec:chapters-ch05:5.4.1}

In summary, it appears that the flagrant fakes have been identified.\footnote[43]{\pinyinterm{li-xueqin} (\citeyear{LiHsuehchin1957Ping}), p. 123, for example, judges \pinyinterm{jingjin} 481 (D), relied on by 陳 \pinyinterm{mengjia-short} (\citeyear{Chen1956Yin}), to be spurious; \pinyinterm{jin-xiangheng} (\citeyear{Chin1962aKufang}) does the same for \pinyinterm{kufang} 1506 (D).} %43
 Skillful
fakes undoubtedly still lurk in the pages of published collections. But by relying as extensively as possible on the inscriptions found in \pinyinterm{jiabian} and \pinyinterm{yibian} and their derivatives,
\pinyinterm{zhuihe} and \pinyinterm{pinbian}, and on those inscriptions which occur repeatedly in collections of
widely differing provenance, the historian can greatly minimize the risk of being deceived
by forgeries. No forger likely to have produced an inscription that will, by itself, significantly distort our understanding of the \pinyinterm{shang}. With these considerations in mind, I
propose that the following scale be used for assessing the reliability of individual oracle
bones.
I.
Attested: applied to bone and shell fragments excavated under scientific conditions
and whose authenticity as \pinyinterm{shang} materials is unquestioned. The \pinyinterm{jiabian}, \pinyinterm{zhuihe}, \pinyinterm{pinbian}, and \pinyinterm{yibian} inscriptions, together with the 1971 and 1973 finds,\footnote[44]{For the 1971 and 1973 finds, see \pinyinterm{guo-moruo-spaced} (\citeyear{Kuo1972Anyang}); KK (\citeyear{KK1975Anyang}), pp. 27--46. On the \pinyinterm{jiabian} finds, see n. 32.} %44
 may be labeled
attested. A fragment from a nonattested collection may attain attested status if it can be
% source: scan 162, printed 145
rejoined, without doubt, to an attested fragment (ideally, the attested fragment should
have been excavated after the nonattested fragment was published, but given the relatively
late publication date of many privately collected inscriptions, that criterion is hard to
satisfy).\footnote[45]{\pinyinterm{kufang} 1661 (D) may be joined with \pinyinterm{jiabian} 297, forming \pinyinterm{zhuihe} 328 (S166.4; 16); \pinyinterm{cuibian} 425 may be joined with \pinyinterm{jiabian} 264, \pinyinterm{jingjin} 1266 with \pinyinterm{pinbian} 24, etc. For other authenticating cases of this sort, see \pinyinterm{zhuixin} 14; 15; 27; 70; 104; 107; see too n. 65. Numerous factors involving bone thickness, fissure marks, and provenance must be considered before a rejoining may be accepted as correct (see sec.~\ref{sec:chapters-ch05:5.7}% was: (sec. 5.7)
).} %45
 On similar grounds, there is a strong presumption that nonattested inscriptions
are also genuine when on the basis of crack numbers, kan-chih date, content, and style,
they appear to belong to the same set, or at least refer to the same event, as an attested
inscription (excavated, if possible, after the publication of the inscriptions in question).\footnote[46]{For example, there is little doubt that \pinyinterm{cuibian} 1108; 1109; \pinyinterm{xubian} 3.11.3; and \pinyinterm{xucun} 1.624 (all S110.2), all from period I, are genuine inscriptions, for they are associated by their cyclical date and wording with such attested inscriptions as \pinyinterm{pinbian} 12.1-2 (sec. 3.7); see \ref{ch:3} (see ch. 3), n. 108. Similarly, \pinyinterm{buci} 632 refers to the same lunar eclipse as \pinyinterm{jiatu} 55 (both S162.1).} %46

2. Acceptable: applied to inscriptions not excavated by archaeologists, but whose
authenticity, on the basis of the criteria suggested in sec.~\ref{sec:chapters-ch05:5.4}% was: (see sec. 5.4)
, there is no reason to doubt.
The bulk of the inscriptions published fall into this category, but no inscription should
be labeled acceptable without explicit review.
3. Uncertain applied to inscriptions which are anomalous in any of the ways
described in sec.~\ref{sec:chapters-ch05:5.4}% was: (see sec. 5.4)
. Uncertain does not necessarily mean inauthentic, but it introduces
a justified note of caution.
4. Spurious: applied to inscriptions which for any of the reasons described above
are thought to be forgeries.
Attention to these categories will encourage historians to assess the authenticity of
their sources. So far as possible, conclusions should be based on attested inscriptions, and
it is for this reason that I list attested inscriptions before nonattested in the footnotes of
this book.
The recognition of previously unidentified fakes, which I have not entered into here,
is not readily reduced to a few simple rules. The basic issue, as with authentication in any
field, is to what extent does the style, placement, vocabulary, and so on of a disputed
inscription accord with or differ from one's educated expectations? Figure 29 reproduces
a series of fake inscriptions, easily recognizable because they are untranslatable; that is,
they violate routine formulas, present nonsense graphs, or juxtapose the graphs in ways
which make no sense and have no \pinyinterm{shang} models. But historians who wish to detect more
subtle forgeries for themselves will learn to do so primarily by reading the relevant scholarship,\footnote[47]{For a general introduction to the problem of forged inscriptions, see Zhang Zongdong (\citeyear{Zhang1970Der}), pp. 30-33. For particular cases, see, for example, \pinyinterm{jin-xiangheng} (\citeyear{Chin1962aKufang}), a particularly instructive case study; \pinyinterm{yan-yiping} (\citeyear{Yen1967Chiaku}); (\citeyear{Yen1974Chia}); \pinyinterm{hu-houxuan} (\citeyear{Hu1973aLintzu}). See too n. 34.} %47
 by studying the large number of fakes that have been identified, and by growing
attuned to the patterns and qualities of genuine inscriptions. Even then, they will rarely
be in a position to identify forgeries unless they actually work with the bones and shells.
45. \pinyinterm{kufang} 1661 (D) may be joined with \pinyinterm{jiatu} 297, forming \pinyinterm{zhuihe} 328 (S166.4;
16); \pinyinterm{cuibian} 425 may be joined with \pinyinterm{jiabian}
264, \pinyinterm{jingjin} 1266 with \pinyinterm{pinbian} 24, etc. For
other authenticating cases of this sort, see \pinyinterm{zhuixin}
14; 15; 27; 70; 104; 107; see too n. 65. Numerous
factors involving bone thickness, fissure marks, and
provenance must be considered before a rejoining
may be accepted as correct (sec.~\ref{sec:chapters-ch05:5.7}% was: (see sec. 5.7)
).
% source: scan 163, printed 146
Authentication, in short, is a matter of connoisseurship and can best be practiced by
handling the actual objects.\footnote[48]{Forgers frequently carved their inscriptions on genuine \pinyinterm{shang} bones (see n. 27). Collagen dating (which would involve destroying the bone) or similar scientific tests might determine the date of the bone, but could tell us nothing of the date of the inscription. On occasion, examination with a loupe or binocular microscope (I have found approximately 16x the most effective magnification) will reveal the fresh traces left by a modern knife, the lack of impacted soil in the groove, and the general cleanness and neatness of the grooves. (For an introductory study of the \pinyinterm{shang} grooves, see \pinyinterm{zhou-hongxiang} [1967/68]; cf. \pinyinterm{zhou-hongxiang} [\citeyear{Chou1976OracleBone}], p. 11.) But the absence of such features is no proof of authenticity. The fake inscriptions in the Musée Guimet, for example (see n. 40) look genuine so far as their incisions are concerned. The bones published as \pinyinterm{bali} 6-8, 10-12 (D) are particularly instructive in this regard, for the same scapula fragments contain fake and genuine inscriptions in close proximity; while it is generally true that the fake grooves appear cleaner and newer and the genuine grooves older and dirtier (as in \pinyinterm{bali} 11), these correlations are frequently relative and subjective. They cannot always be used as authenticating criteria; other criteria, such as writing style and the sense or nonsense of the inscription must still be considered. See too Barnard's comments cited at the start of sec.~\ref{sec:chapters-ch05:5.4}% was: (see sec. 5.4)
.}

\section[Methods of Reproduction]{Methods of Reproduction}
\label{sec:chapters-ch05:5.5}

No scholar is in a position to examine every bone and shell cited. Many are
scattered over the globe; others have been lost or at least have been lost track of.\footnote[49]{As in the case of \pinyinterm{minghou}; we have Menzies' rubbings, but the whereabouts of the 2,812 fragments is unknown (\pinyinterm{minghou} preface, p. 3). For other cases in which the original oracle bones have now been lost, see Lefeuvre (\citeyear{Lefeuvre1975Les}), pp. 27-28, n. 1, pp. 41, 49, n. 3. For notices about fragments disappearing and reappearing in collections in the United States, see \pinyinterm{zhou-hongxiang} (\citeyear{Chou1976OracleBone}), pp. 16, 18, 22, 23.} For
most research, therefore, the scholar will rely upon three kinds of reproduction: rubbings,
drawings, and photographs.\footnote[50]{The present corpus consists of about 43 collections of rubbings (approximately 38,000 pieces), about 20 collections of drawings (12,654 pieces), and various small collections of photographs (probably less than 1,000 pieces). These estimates are based on \pinyinterm{dong-short} (\citeyear{Tung1953Chia}), pp. 375-378, adding Menzies, \pinyinterm{minghou}, US, and White. A rubbing reproduces the inscription mechanically; but unlike printing, where the ink lies between the object and the paper and a negative image is reproduced, the paper lies between the object and the ink and a positive image is reproduced. The term ``squeeze'' is apparently preferred in Western classical scholarship. For the general technique, see Woodhead (\citeyear{Woodhead1967Study}), pp. 78-83; for details about making oracle-bone rubbings, see \pinyinterm{zhou-hongxiang} (\citeyear{Chou1976OracleBone}), pp. 5, 18, n. 15. It may be remarked that, as modern rubbings attest, the carved scapulas and plastrons were potential ``printing blocks.'' References to smearing the plastrons with ink do appear in \pinyinterm{zhou-dynasty} texts (\ref{ch:2} (see ch. 2), n. 155), but there is no evidence that the \pinyinterm{shang} reproduced copies of their divinations in this way. The problem of why the graphs were carved will be considered in Studies.} Usually, the scholar working with a particular inscription
must be content with one or, at best, two forms of reproduction.\footnote[51]{For an example of the three types of reproduction, consider \pinyinterm{jiabian} 2122. A photograph of the actual plastron is given in \pinyinterm{dong-short} (\citeyear{Tung1931Takuei}), plastron 4; (\citeyear{Tung1948aTen}), no. 6; and \pinyinterm{tongzuan}, ``pieh'' 2, IA; \pinyinterm{jiabian} 2122 and 家土 80 are rubbings of the plastron; a drawing is found in \pinyinterm{dong-short} (\citeyear{Tung1945Yin}), pt. 2, ch. 5, table 3, p. 10b. For a rubbing, photograph, and drawing of the same fragment (USS 651) juxtaposed, see \pinyinterm{zhou-hongxiang} (\citeyear{Chou1976OracleBone}), p. 5.}
Rubbings are generally more reliable than drawings, and more reliable, in some
respects, than photographs, though the way in which rubbings were made and printed
% source: scan 164, printed 146
may affect their usefulness. 52 Photographs\footnote[52]{Not infrequently a rubbing may have been reproduced in two or more collections. and one rubbing may be clearer than the other; e.g., \pinyinterm{waibian} 453 is a far clearer rubbing than \pinyinterm{xucun} 1.632. (For an error caused by an early unclear rubbing of \pinyinterm{jiabian} 1289, see \pinyinterm{dong-short} [\citeyear{Tung1952Putzu}], p. 284). \booktitle{Tiexin} reprints the best rubbings made of the \booktitle{Tieyun} pieces; the quality of reproduction in the original edition was not always good. Similarly, \pinyinterm{jicheng} 1 initiates a major attempt to reproduce mint rubbings from various collections and to identify duplicates; until an index is provided, however, it will be hard to use. (For a useful review of \booktitle{Tiexin} and \pinyinterm{jicheng} 1 see Mickel [\citeyear{Mickel1976bThree}].) \pinyinterm{shima-short}'s \booktitle{Sōrui} (sec. 3.3.2) lists some examples of duplicate rubbings, but a comprehensive catalog is needed. The work of Mickel (\citeyear{Mickel1974Reduplicated}), (\citeyear{Mickel1974aIndex}), (\citeyear{Mickel1976aIndex}), and (\citeyear{Mickel1977aIndex}) helps satisfy this need. With regard to the quality of editions, Jung Keng (\citeyear{Jung1947Chiaku}), p. 22, notes, for example, that the first collotype edition of \booktitle{Shiyi} was not at all clear and that the lithographic version was still worse. For a detailed discussion of the various methods of printing (collotype, metal plate, rubber plate, photolithography, etc.) and types of rubbing (``glossy dark'' or ``cicada wing''), and examples of each, see \pinyinterm{chen-mengjia-spaced} (\citeyear{Chen1956Yin}), p. 48; \pinyinterm{dong-short} (\citeyear{Tung1958Chiaku}), pp. 918-920; \pinyinterm{zhou-hongxiang} (\citeyear{Chou1976OracleBone}), pp. 2-7, 18. The quality of the various editions of rubbing collections will be fully treated in Lefeuvre's forthcoming work (see \ref{ch:3} (see ch. 3), n. 4, above). It is occasionally necessary to beware rubbings that have been spuriously improved to produce graphs that are neater and clearer than those on the original oracle bone; see, for example, the refurbished rubbings of Wang Ziyu (\citeyear{Wang1933Chia}); see too the critical comments of \pinyinterm{sun-haibo} (\citeyear{Sun1937Fushi}) on the way the \pinyinterm{fuyin} rubbings were ``improved."} have one clear advantage over rubbings in their
ability to reproduce graphs that were brush-written but not inscribed;\footnote[53]{See \ref{ch:2} (see ch. 2), n. 96.}5³ their accuracy,
if unretouched, and the speed with which they may be made also recommend them.
Photographs, nevertheless, have their drawbacks; the graphs frequently offer insufficient
contrast with the bone surface to be photographed with success; it is not always possible
to bring all the curved surface of a bone into sharp focus; erasures and corrections, which
may appear in rubbings, do not photograph well.\footnote[54]{On the difficulties involved, see, for example, Britton (\citeyear{Britton1935Yin}), pp. 1, 13; \pinyinterm{li-yan-spaced} (\citeyear{LiYen1970PeiMei}); \pinyinterm{zhou-hongxiang} (\citeyear{Chou1976OracleBone}), pp. 3, 5, 22. The virtues of rubbings and photographs may readily be compared by consulting \pinyinterm{haiwai}, pt. 1, 1–12, which gives both; see too \pinyinterm{zhou-hongxiang} (\citeyear{Chou1976OracleBone}), pp. 2-7. It should be noted that pu cracks, especially when not engraved (sec.~\ref{sec:chapters-ch02:2.10}% was: (sec. 2.10)
), are frequently too thin or subtle to appear in a rubbing (e.g., the cracks in Menzies 2455; 2528; 2664 but are reproduced well by photographs; on this point, cf. \pinyinterm{zhang-bingquan} (\citeyear{Zhang1956Pu}), p. 241.} %54
 Drawings, in particular, must be
treated with caution, especially when the draftsman was unfamiliar with the content of
the inscriptions.\footnote[55]{奇-chi, \pinyinterm{jingjin}, \pinyinterm{xucun}, pt. 2, \pinyinterm{nanbei}, \pinyinterm{ninghu}, and 殷墟 reproduce only drawings of inscriptions. Menzies' drawings in 殷墟 are not of high quality (Jung Keng [\citeyear{Jung1947Chiaku}], p. 21) and those of Chalfant (in 奇-chi, \pinyinterm{jingjin}, and \pinyinterm{kufang}) are thought less good (ibid., pp. 19, 2013 for examples of Chalfant's errors, see \pinyinterm{zhou-hongxiang} (\citeyear{Chou1976OracleBone}), pp. 2, 3. On the way Chalfant's hand improved with practice, see 奇-chi, preface, p. 1. \pinyinterm{hu-houxuan}'s copies in \pinyinterm{nanbei} and \pinyinterm{ninghu} contain many errors (\booktitle{Shitao}, pt. 2, preface, p. 3a; see too \pinyinterm{dong-short} and \pinyinterm{jin-xiangheng} [\citeyear{TungChin1961Penhsi}], p. 72). One of the master draftsmen was \pinyinterm{dong-zuobin}; see, for example, his numerous drawings in 殷曆譜 (\pinyinterm{dong-short} [\citeyear{Tung1945Yin}]); for his drawings of \pinyinterm{yibian} 107643, see \pinyinterm{zhongguo-wenzi}, nos. 2-3, 5-13, 16-17 (1961-1965). In all these cases, it is salutary to bear in mind a comment by Creel ([\citeyear{Creel1938Studies}], p. 108): ``It has been proved time after time that when one tries to draw these [oracle-bone] characters in facsimile it is virtually impossible to prevent his mental preconceptions from influencing the copy, try as he will to make nothing more nor less than a simple reproduction.'' For the way different drawing styles may influence judgments about \pinyinterm{shang} calligraphy, compare \pinyinterm{xu-jinxiong}'s drawings (Menzies, plates 265-281) with those of \pinyinterm{yan-yiping} (\pinyinterm{zhuixin} 566-594) for the same pieces. Note in this regard, however, that rubbings are not infallible either; thickness of character stroke, for example, may vary with ``the kind of ink and paper used and the pressure applied'' by the rubbing maker (\pinyinterm{zhou-hongxiang} [\citeyear{Chou1973Computer}], p. 181). 147 ]} %55
 Among other factors, graphs may be wrongly drawn, cracks, crack
% source: scan 165, printed 147
numbers, and crack notations omitted, and failure to reproduce plastron fissures may
handicap attempts at reconstruction. The fact that an inscription is a hand copy should
always be noted in the gloss or citation (done in this book with the letter D). The current
trend to reproduce rubbings (or photographs) of collections which had formerly been
reproduced only as drawings is greatly welcomed.\footnote[56]{E.g., \pinyinterm{minghou}; 台-ta 2; US; \pinyinterm{jin-xiangheng} (\citeyear{Chin1973Fujen}); \pinyinterm{hu-houxuan} (\citeyear{Hu1973aLintzu}). Photographs of 68 pieces originally reproduced as drawings in 奇-chi, W, have been published in \pinyinterm{haiwai}.} %56

Certain caveats should be applied to all forms of reproduction: one cannot always tell
if the inscription is on bone or shell, one cannot tell its thickness (an important consideration
in rejoining fragments; sec.~\ref{sec:chapters-ch05:5.7}% was: (sec. 5.7)
), one cannot tell what hollow forms or even graphs may
have been on the back, one cannot even be sure that the shape of the entire fragment has
been reproduced.\footnote[57]{The top of \pinyinterm{yicun} 888, for example, was cut off; \pinyinterm{hu-houxuan} (\citeyear{Hu1964Chiakuwen}), p. 138 and fig. 3, restores it. Other rubbings were cut in two; e.g., \pinyinterm{fuyin}, ``Ti'' 207 and ``Wen'' 40 (see the complete fragment at \pinyinterm{xubian} 1.45.3 or \pinyinterm{yicun} 914 [胡, op. cit., p. 156]); \pinyinterm{qianbian} 1.14.7 and 1.16.3 (rejoined at \pinyinterm{chaocun} 11). Even the \pinyinterm{yibian} rubbings do not always reproduce the whole fragment; e.g., the upper left quadrant of fragment F of \pinyinterm{pinbian} 395 is not reproduced in \pinyinterm{yibian} 3063. In his discussion of the oblong tablets into which oracle bones were occasionally sawn (sec.~\ref{sec:chapters-ch01:1.3.1}% was: (sec. 1.3.1)
), \pinyinterm{liu-yuanlin} (\citeyear{Liu1969Kuchien}), pp. 237–243, notes that the existence of such tablets had long been unrecognized because early books of rubbings did not always reproduce the full shape of the bone or shell fragment.} %57
 It goes without saying that scholars should work with the actual bones
or shells when possible.\footnote[58]{It is probable that my discussion of \pinyinterm{pinbian} 12-21 (sec.~\ref{sec:chapters-ch03:3.7}% was: (sec. 3.7)
) would require revisions if I had the opportunity to examine the plastrons themselves. For examples of the dangers involved when the actual oracle bone is not studied, see \pinyinterm{yan-yiping} (\citeyear{Yen1961Chiaku}), pp. 516-520.} %58

Methods of Transcription
The translation by transcription of oracle bone graphs into modern Chinese
graphs introduces (as already suggested (sec.~\ref{sec:chapters-ch03:3.5}% was: (sec. 3.5)
)) the risk of distortion. It is a classic case
of traduttore, traditore, all the more insidious because likely to go unremarked.\footnote[59]{On some of the problems involved, see Serruys (\citeyear{Serruys1974The}), pp. 15-19, and, for a particularly instructive case of erroneous transcription, p. 101, n. 23.} %59
 Ideally,
every inscription quoted should be presented as a rubbing and be accompanied by a tracing
of the \pinyinterm{shang} characters, followed by a transcription into modern characters and then by
a translation into a modern language. Such a method, in a work which cites several hundred
inscriptions, would be impossibly cumbersome. Some scholars have chosen to write all
the inscriptions they discuss in \pinyinterm{shang} form alone.\footnote[60]{E.g., \pinyinterm{shima-short} (\citeyear{Shima1958Inkyo}); Zhang Zongdong (\citeyear{Zhang1970Der}); Takashima (\citeyear{Takashima1973Negatives}).} %60
 I have followed their purist example
by providing in the margin \pinyinterm{shima-short}'s oracle-bone transcriptions, taken from \booktitle{Sōrui} (sec.~\ref{sec:chapters-ch03:3.3.2}% was: (sec. 3.3.2)
),
for some of the inscriptions translated in this book. For the scholar, steeped in \pinyintext{jiaguwen},
such transcriptions will generally suffice, at least as transcriptions (see the next paragraph).
% source: scan 166, printed 148
But the familiar, modern forms are needed if the inscriptions and the scholarship about
them are to be read by nonspecialists. As a general rule, I would suggest that the modern
graph form may be used when inscriptions are cited for purposes of illustration and no
issue involving the identity of the word or the form of the graph is being raised. When the
word or the form of the graph is in question, it is essential that the chia-ku graph itself be
provided and that the modern transcription, if one is used (as it is usually convenient to
do in the case of names, for example), preserve contact with the ancient epigraphy by
conforming to the structure of the \pinyinterm{shang} graph, rather than merely providing a modern
character equivalent.\footnote[61]{This means, for example, that (S477.3) of (S12.4) as ssu E, rather than ch'iu , or of should be transcribed as \#, not as tso or ``F, or (S303.1) as huo involves much scholarship even as . By contrast, the transcription of (provided, in these cases, by CKWT, pp. 93-95, (figs. 12 and 13) as mien may express the presumed meaning; it does not reproduce the structural components of the original graph. The importance of such faithful reproduction is convincingly argued by Barnard, e.g., (1958), p. 15, n. 4; (1959a), p. 8.} %61
 If this is not done, important clues may be lost in the translation.
For graphs where no modern equivalent can be proposed, the \pinyinterm{shang} form must be used.
This means that some inscriptions may be presented in mongrel style, eight parts modern,
one part modified modern, one part ancient. This is inconsistent, but it is the most efficient
approach for a work of history, and there is ample precedent.
All transcriptions are to be taken as mere representations of the \pinyinterm{shang} document.
They are smooth, and hence potentially false, depictions of what may have been a wild
and rugged epigraphic landscape.\footnote[62]{To cite but two examples: the transcription as tsu or as pit (cf. \ref{ch:3} (see ch. 3), n. 47, n. 51) supplies radicals and thus amends the \pinyinterm{shang} graphs.} %62
 All study of the inscriptions must finally be based on
the original graphs in their original context. It must be based on the rubbings and, if
possible, on the original scapulas and plastrons. As a matter of convenience, I have given,
when it is useful to do so, the Sōrui reference after each inscription cited; this will enable
the reader without access to the rubbing collections to see what the \pinyinterm{shang} graphs looked
like in each case. But \pinyinterm{shima-short} presents a standardized form of oracle-bone script; he has
not always copied the inscriptions accurately or consistently; and he gives no indication
of placement and direction on the bone. His entries too must be treated as representations.
Oracle-bone scholars use certain standard conventions in their transcriptions. The
symbol is used when we cannot tell if a character is missing, or, if so, how many characters
have been lost. The symbol is used to represent one missing character.\footnote[63]{E.g., \pinyinterm{jiabian} kǎoshì, ``凡例,'' p. 9.} %63
 Parentheses
indicate that meaning has been supplied. Square brackets indicate that the graphs are
not present on the oracle bone we have, but are thought to have been present originally
and have been restored. The symbols \_ or \_, inserted in some modern transcriptions
(such as \pinyinterm{pinbian} kǎoshì), indicate that the text columns were written downward and
to the left or downward and to the right, respectively.
% source: scan 167, printed 149
Reconstruction of Inscriptions
Many of the oracle bones excavated at \pinyinterm{xiaotun} were already in fragments when
they came out of the ground; others were broken later by accident or to be sold separately.
Rejoining these fragments, and thus the inscriptions, is a work of vital importance.\footnote[64]{On the method of rejoining fragments, see \pinyinterm{dong-short} (\citeyear{Tung1945Yin}), pt. 2, ch. 5, pp. 22a ff.; ch. 8, pp. 12a-13a; (\citeyear{Tung1954aYinli}), p. 120; (\citeyear{Tung1958Chiaku}), pp. 915-918; \pinyinterm{zhang-bingquan} (\citeyear{Zhang1956aPutzu}), pp. 242-243; Serruys (\citeyear{Serruys1974The}), pp. 15-19. Woodhead (\citeyear{Woodhead1967Study}), pp. 67-76, is also useful for the \pinyinterm{shang} scholar. Note that a small part of the rejoiner's work may involve rejoining fragments which were once whole, or rubbings which were subsequently cut in pieces. For example, \booktitle{Tieyun} 108.1, originally published as one fragment, was subsequently broken in two and published as \pinyinterm{xucun} 1.358 and \booktitle{Shitao} 2.92 (see \pinyinterm{zhuixin} 285). For rubbings, see n. 57.} %64
 It
permits us to reconstitute whole divination sentences, to place them in the context of
adjacent inscriptions, and even, on occasion, in the context of a series of inscriptions.
Rejoinings may also serve, as already noted (sec.~\ref{sec:chapters-ch05:5.4.1}% was: (sec. 5.4.1)
, par. 1), to authenticate an inscription
that was not scientifically excavated.\footnote[65]{See n. 45. The fact that \pinyinterm{jingjin} 1266 may be joined to \pinyinterm{pinbian} 24 (see \pinyinterm{pinbian} 24, kǎoshì, pp. 47-51) suggests either that the excavators of YH127 were careless and the piece was stolen or that the plastron was broken before the \pinyinterm{shang} put it in the ground. Cf. \pinyinterm{yan-yiping} (\citeyear{Yen1973Kuanyu}), pp. 2b, 6a, 6b-7a, for this and two similar cases.} %65
 Certain collections are composed entirely of rejoined
fragments; and many other joinings have been made on an individual basis.\footnote[66]{The major published collections of rejoined pieces are \pinyinterm{chaocun}, \pinyinterm{zhuihe}, \pinyinterm{hebian}, and \pinyinterm{pinbian}; \pinyinterm{jiatu} (plates 1-58 at the rear of \pinyinterm{jiabian} kǎoshì) contains 221 rejoinings of \pinyinterm{jiabian} fragments. \pinyinterm{zhuixin} reproduces by rubbing and drawing many of these rejoinings (with the exception of those involving \pinyinterm{yibian} fragments) and many new ones; the preparation of an index would render this useful work more valuable. For a review of \pinyinterm{zhuixin} see Mickel (\citeyear{Mickel1976bThree}). Mickel (\citeyear{Mickel1974Reduplicated}), (\citeyear{Mickel1974aIndex}), (\citeyear{Mickel1976aIndex}) and (\citeyear{Mickel1977aIndex}) provide indexes to the rejoined (and reduplicated) fragments in \pinyinterm{jianshou}, \pinyinterm{qianbian}, \pinyinterm{jinghua}, \pinyinterm{fujia-title}, \pinyinterm{fuyin}, \pinyinterm{houbian}, \pinyinterm{xubian}, Kikkō, \pinyinterm{buci}, \booktitle{Tieyi}, \booktitle{Tieyu}, \booktitle{Tieyun}, \pinyinterm{cunzhen}, \pinyinterm{dong-short} (\citeyear{Tung1929aHsin}), \pinyinterm{yicun}, and 殷墟 (see table 9).} %66
 The work
of rejoining, which may be likened to an attempt to put together hundreds of separate,
yet deceptively similar, jigsaw puzzles from thousands of similar pieces, many of which
are missing, some of which have lost their original shape (and some of which may even
be rejoined in more than one way),'' may eventually be speeded up by the application
of computer technology.\footnote[67]{For the correct rearrangement of the fragments wrongly rejoined in \pinyinterm{jiatu} 5, see n. 76.} %68
\footnote[68]{\pinyinterm{zhou-hongxiang} (\citeyear{Chou1973Computer}).}
\pinyinterm{dong-zuobin}'s tongue-in-cheek claim that if one knows the kan-chih date, inscription
type, and structure of the bone and shell, reconstruction is ``easier than putting a broken
mirror together''\footnote[69]{\pinyinterm{dong-short} (\citeyear{Tung1954aYinli}), p. 120.} should remind us that any rejoining, especially if made only on the
basis of rubbings, must be treated with caution. The rejoiner needs to be aware of the
following pitfalls: rubbings of individual fragments may stretch or shrink to different
degrees so that one rubbing will not match another, even though the original fragments
do; it may be difficult to fit rubbings together since the curvature of the reconstituted
bone or shell will differ from that of the rubbings made from individual fragments; differing
pressures and conditions underground may have warped, broken, or discolored fragments
from the same piece in different ways so that even the fragments themselves will not fit
% source: scan 168, printed 150
well or appear related; impacted dust or dirt may prevent suture lines from appearing
in one or more of the rubbings; and rubbings may be of such varying quality that a rejoining
may, at first glance, appear out of the question.\footnote[70]{\pinyinterm{pinbian} 24, kǎoshì, pp. 47-48; \pinyinterm{zhou-hongxiang} (1970. pp. 54-56; (\citeyear{Chou1973Computer}), p. 180.} In addition, it is essential that the rejoiner
have a thorough understanding of the morphology of scapulas, plastrons, and carapaces,
as well as a firm grasp of the way hollows, cracks, crack numbers, crack notations, and
inscriptions were likely to be placed on different sections of bone or shell (secs.~\ref{sec:chapters-ch01:1.5.1}; \ref{sec:chapters-ch02:2.4};
\ref{sec:chapters-ch02:2.6}; \ref{sec:chapters-ch02:2.9.4}% was: (secs. 1.5.1; 2.4; 2.6; 2.9.4)
).\footnote[71]{On this aspect of reconstruction, see \pinyinterm{dong-short} (\citeyear{Tung1945Yin}), pt. 2, \ref{ch:5} (see ch. 5). pp. 19a-21b; ch. 9, p. 45a; \pinyinterm{zhang-bingquan} (\citeyear{Zhang1956Pu}), pp. 240, 243; \pinyinterm{yan-yiping} (\citeyear{Yen1967Chiaku}), pp. 1 6-12.} \pinyinterm{dong-short} himself made some brilliant reconstructions, fascinating examples of
epigraphic detective work (for example, fig. 30),\footnote[72]{E.g., \pinyinterm{dong-short} (\citeyear{Tung1945Yin}), pt. 2, \ref{ch:3} (see ch. 3), p. 32b; \ref{ch:5} (see ch. 5), pp. 19a-21b; ch. 6, pp. 6b-8a; ch. 7, p. 2a; ch. 9, pp. 54a-58b: ch. 10, pp. 5b-6a. Many, if not all, of these rejoinings are reproduced by \pinyinterm{zhuixin} 296-348.} but some are speculative, and others he
subsequently admitted to be wrong.\footnote[73]{\pinyinterm{dong-short} (1948b, pp. 186-188; (\citeyear{Tung1952Putzu}), p. 282.} A fair number of erroneous reconstructions have
been published, a few with serious consequences.\footnote[74]{For references to erroneous reconstructions, see \pinyinterm{dong-short} (\citeyear{Tung1945Yin}), pt. 2, ch. 7, p. 3b; \pinyinterm{yan-yiping} (\citeyear{Yen1954Payueh}); \pinyinterm{li-xueqin} (\citeyear{LiHsuehchin1956Tan}), p. 16; \pinyinterm{qu-wanli} (\citeyear{Chu1965aShihchi}), pp. 73-74. \pinyinterm{zhang-bingquan} (\citeyear{Zhang1956aPutzu}), p. 244, joins \pinyinterm{yibian} 4169 to \pinyinterm{yibian} 2564+3139; but \pinyinterm{yibian} 4169 does not appear in \pinyinterm{pinbian} 477; the original rejoining was apparently rejected. Similarly, \pinyinterm{yan-yiping} (\citeyear{Yen1961Chiaku}), 77 and 235, combines \pinyinterm{yibian} fragments not found in \pinyinterm{pinbian}. \pinyinterm{zhuixin}, ch. 10, passim, challenges ninety-seven rejoinings made by previous scholars; \pinyinterm{yan-yiping}'s discussions of such points as differing edge angle, wrong reconstruction of the character straddling the break, and differences in writing style are instructive; she also notes wrong rejoinings in the body of the work, as at \pinyinterm{zhuixin} 22, 90, 246, 266, 314, 337, 470, and 471 (the last two concern the faulty linking of front and back rubbings which do not, in fact, reproduce the front and back of the same oracle bone). Wrong or incomplete reconstructions proved particularly misleading in the attempt to date the lunar eclipse records; for the revised rejoinings, see \pinyinterm{dong-short} (\citeyear{Tung1952Putzu}), pp. 282-286; \pinyinterm{zhang-bingquan} (\citeyear{Zhang1956aPutzu}), pp. 244–245. These and other eclipse records will be discussed in Studies.} It goes without saying that further
fragments may be added to published rejoinings until the whole plastron or scapula is
completed.\footnote[75]{For example, \pinyinterm{zhuixin} 4, 14, 16, etc. supplement some of the \pinyinterm{jiatu} rejoinings. As \pinyinterm{yan-yiping} notes (\pinyinterm{zhuixin}, preface, p. 1b), not all the \pinyinterm{yibian} rejoinings have yet been made.}
In conclusion, the historian should be cautious in using such restorations unless the
scholar who reassembled the pieces was able to work with the oracle bones themselves and
note, in particular, that thickness, fissure lines, and other physical features matched.\footnote[76]{Even working with the fragments themselves does not guarantee a correct rejoining. \pinyinterm{yan-yiping} (1962 [cf. 1961, 286]) by using rubbings was able to correct the faulty rejoining of \pinyinterm{jiatu} 5, which had been based on the fragments.}
And even in those cases, the match can only be thoroughly reliable if an inscription (and
preferably a graph or series of graphs) straddles the break between the rejoined fragments.
Certain reconstructions are of the ``exploded'' variety; that is, they identify and
place fragments on the same scapula or plastron, even though the intervening fragments
are missing; and they may further supply missing graphs on the basis of context and expecta-
tion (for example, fig. 30).\footnote[77]{E.g., \pinyinterm{dong-short} (\citeyear{Tung1945Yin}), pt. 2, ch. 5, pp. 19a-21b; ch. 6, p. 7b; ch. 7, p. 2a, etc.; \pinyinterm{zhuixin} 72-75; 331 363-365; 479; 546, etc.} Similarly, missing graphs may be supplied for single fragments
when part of the inscription has been broken off (as in sec.~\ref{sec:chapters-ch03:3.6.1}% was: (sec. 3.6.1)
; fig. 20). In all cases,
% source: scan 169, printed 151
\pinyinterm{shima-short}'s \booktitle{Sōrui} (sec.~\ref{sec:chapters-ch03:3.3.2}% was: (sec. 3.3.2)
) is invaluable in giving the inscription patterns for any particular
topic; in particular, the regularity of the ritual cycle in periods IIb and V assists the restoration of inscriptions from those periods.\footnote[78]{For the ritual cycle see \ref{ch:4} (see ch. 4), n. 102; appendix \ref{app:5}, sec. \ref{sec:appendices-app05:2}% was: (appendix 5, sec. 2)
.} %78
 Such reconstructions, provided the graphs supplied
are always enclosed in square brackets and the reasons for the reconstruction given, when
necessary, in a gloss, have an important role to play. Many-such as the restoration of a
cyclical day date or diviner's name-are routine, and appear frequently in inscription
commentaries. Others may involve far greater scope. Exploded and imaginative reconstructions should be treated with even more caution than that required for those rejoinings
in which the fragments are contiguous.\footnote[79]{Several of the erroneous reconstructions cited in n. 73 and n. 74 were of this exploded or imaginative variety.} %79

\section[Grounds for Cautious Optimism]{Grounds for Cautious Optimism}
\label{sec:chapters-ch05:5.8}
The oracle-bone inscriptions, documents of a kind unique in history, provide a
vivid, if sporadic, sense of contact with the Chinese Bronze Age. The highly polished
scapula fragments, with the look and feel of old carved ivory, may be called beautiful.
Sanctified by their great age, eloquent witnesses to a royal religion and belief, they invite
reverent handling. Those who spent years in the heroic yet depressing task of shuffling
and sorting thousands of oracle-bone fragments, may take a less romantic view of the
matter.\footnote[80]{See, e.g., \pinyinterm{zhang-bingquan}'s description of the labors involved in rejoining the \pinyinterm{yibian} fragments (\pinyinterm{pinbian}, preface, p. 1).} %80
 Nevertheless, the occasional chance to handle an oracle bone of museum quality
affords a pleasure, a sense of immediate contact with ancient and holy things, denied to
historians working, say, with Sung recensions of \pinyinterm{zhou-dynasty} or \pinyinterm{han-dynasty} documents. The oracle bones
were the very bones and shells that the \pinyinterm{shang} kings and their diviners cracked and prognosticated over three millennia ago.
Like all documents, the inscriptions have their historiographical advantages and
disadvantages. The following points may be made in their favor: \listitem{1} The divination
inscriptions are indisputably old (for the question of forgeries, sec.~\ref{sec:chapters-ch05:5.4}% was: (see sec. 5.4)
). \listitem{2} The original
inscriptions, which lend themselves so conveniently to the making of modern rubbings,
may be reproduced without the risk of scribal or typographical error; rubbings are copies
made from mechanical contact with the originals; they are not copies of copies. \listitem{3} The
inscriptions were lost for some three thousand years; they have thus reached us undistorted
by errors in transmission and unaccompanied by tendentious commentaries (though
modern scholars, of course, have not lacked axes to grind in interpreting the inscriptions).
\listitem{4} They are not epiphenomenal or decorative documents; they are true records of how
the \pinyinterm{shang} king actually conceived of his alternatives. \listitem{5} They have been preserved by
% source: scan 170, printed 152
accident, without any evident intent to influence posterity.\footnote[81]{Preliminary research, which will be presented in Studies, suggests that the oracle bones were not engraved or stored for posterity; they come to us from the cellars and refuse pits of the \pinyinterm{shang}, not from the temple archives.} %81
 The oracle-bone inscriptions,
in short, are ``hard'' documents in the best historiographical sense. This is not to say that
they were documents with no interest to serve; if that were so they would not have been
recorded. They have their biases, but they are the genuine and contemporary biases of
the recorders and not the ideological biases of later transmitters;\footnote[82]{The preference for recording only auspicious crack notations and validating verifications has already been noted (secs.~\ref{sec:chapters-ch02:2.6}, \ref{sec:chapters-ch02:2.8}% was: (secs. 2.6, 2.8)
).} %82
 the texts say what their
authors intended they should say. The divination inscriptions have the invaluable quality
of setz im leben, immediately apprehended by the recorder and accurately transmitted to us.
On the other hand, there are disadvantages: \listitem{1} The oracle-bone inscriptions may
have been transmitted accurately, but it is frequently hard, and occasionally impossible,
to determine accurately the meaning of what has been transmitted sec.~\ref{sec:chapters-ch03:3.5}% was: (see sec. 3.5)
). \listitem{2} Further,
the inscriptions present a narrow view. They deal primarily with ritual and sacrifice and
do so in standard forms; they deal exclusively with the concerns of a small elite group;
they tell us little, for instance, of economic organization, or of the life of the peasantry.
\listitem{3} We have only a small sample of the total number of inscriptions that were made during
the last eight or nine reigns of a dynasty that supposedly saw some thirty kings on the
throne. For these reasons, \pinyinterm{dong-zuobin} concluded that the inscriptions were not representative, that they could not be used to understand more than 1 percent of \pinyinterm{shang} culture,
that and here he was addressing an international symposium-they represented ``only
evaporated and left little
the dregs of a civilization, the best parts of which have.
trace.''\footnote[83]{\pinyinterm{dong-short} (\citeyear{Tung1949Yinhsu}), pp. 241-242; (\citeyear{Tung1951cChinese}), p. 5; (\citeyear{Tung1951dChungkuo}), p. 28. For his more optimistic summary of what the oracle bones can tell us, see \pinyinterm{dong-short} (\citeyear{Tung1951Lun}); cf. \pinyinterm{li-xueqin-spaced} (\citeyear{LiHsuehchin1959Kuan}), p. 22.} And H. G. Creel finds ``the state of our knowledge of the \pinyinterm{shang} period ...
exasperating. We know so much about a few things that it seems intolerable that we know
so little about the culture and history as a whole.\footnote[84]{Creel (\citeyear{Creel1970Origins}), p. 31. Cf. Van Gulik (\citeyear{Van1961Sexual}), p. 4: ``We find ourselves in approximately the same position as a man who in say the year 5000 would try to reconstruct our present Western civilization while having at his disposal only a collection of tombstones from a number of cemeteries now scattered over Europe."}''
``84
Such pessimism is not to be dismissed, but it can be qualified. We may be in the
position of somebody attempting to reconstruct classical Greek culture by using the
archives of Delphi; clearly our ``diviner's eye view'' will be narrow and special. But the
attempt must be made if we are to have more than an antiquarian's knowledge of the
origins of Chinese civilization. And it may be made, I believe, with some optimism.
\pinyinterm{xiaotun} was not Delphi; it was a major, and probably the major, political and religious
center of the North China plain.\footnote[85]{On the status of \pinyinterm{xiaotun}, see \ref{ch:2} (see ch. 2), n. 2. 153 == 1} %85
 The \pinyinterm{shang} king was no mere Pythia or priest; he was
a powerful theocrat presiding over economic, military, and religious structures of considerable administrative complexity. The inscribed charges and prognostications did
% source: scan 171, printed 153
not vary in coherence (as they did at Delphi),\footnote[86]{Park and Wormell (\citeyear{ParkeWormell1956Delphic}), p. 33.} %86
 requiring ambiguous explication by a
priest; they were recorded in clear. Finally, most of the divinations deal not with the
extraordinary but with the ordinary, not with emergencies but with the routine business
of \pinyinterm{shang} religion and government as it was practiced on a day-to-day basis. Ideally we
would like to have far more evidence; but we can make good and full use of what we have.
Initial studies have already been made, in Western languages alone, of such topics
as ancestor worship, the role of the Nature powers, conceptions of disease or disaster,
kinship organization, succession and descent, the \pinyinterm{shang} calendar, administrative organization, the genesis of urban forms, public work, and legitimation.\footnote[87]{See, for example, \pinyinterm{zhang-guangzhi} (\citeyear{Chang1964Some}); (\citeyear{Chang1974Urbanism}); \pinyinterm{zhou-hongxiang} (\citeyear{Chou1968SomeAspects}); \textcite{Keightley1969Public}; (\citeyear{Keightley1975The}); (\citeyear{Keightley1975Legitimation}); Zhang Zongdong (\citeyear{Zhang1970Der}); Zhou Lin (\citeyear{Chao1970Marriage}); (\citeyear{Chao1972ShangGovernment}); Wheatley (\citeyear{Wheatley1971Pivot}); Vandermeersch (\citeyear{Vandermeersch1975Wangdao}); \textcite{Mickel1976Semantic}; Nivison (\citeyear{Nivison1977Interpretation}). Recent studies in Chinese and Japanese have used the inscriptions to study, for example, \pinyinterm{shang} agriculture (\pinyinterm{yu-xingwu} [\citeyear{Yu1972Tsung}]), kinship organization (\pinyinterm{zhang-guangzhi} [\citeyear{Chang1973Tan}]), metallurgy (Yan Yun [\citeyear{YenYun1973Shangtai}]), punishments (\pinyinterm{hu-houxuan} [\citeyear{Hu1973Yintai}]), sericulture (\pinyinterm{hu-houxuan} [\citeyear{Hu1972Yintai}]), technology (\pinyinterm{li-shi} [\citeyear{LiShih1976Yu}]), slavery (\pinyinterm{hu-houxuan} [\citeyear{Hu1976Chiakuwen}]; Yang and Yang [\citeyear{Yang1977Tsung}]), and religion (Akatsuka [\citeyear{Akatsuka1977Chugoku}]).} %87
 Much of what we
already know about the \pinyinterm{shang}, including the fact that the dynasty itself existed, we owe
to the oracle-bone inscriptions. Without them we would see the pre-\pinyinterm{zhou-dynasty} Bronze Age
of North China far more dimly than we do.
Pessimism about the contribution that the oracle inscriptions can make has, I believe,
been due in large part to the attempt to fit the \pinyinterm{shang} into the standard framework of a
Chinese dynastic history. We cannot write annals of the \pinyinterm{shang} emperors, biographies of
ministers and generals, treatises on economics or water conservancy. It would be folly
to believe, barring new and revolutionary discoveries excavated from the soil of China,
that we will ever be able to write of the \pinyinterm{shang} in that way.\footnote[88]{Hughes (\citeyear{Hughes1961Consciousness}), p. 9.}
``88
Students of \pinyinterm{shang} history may, and indeed must, attempt to work as political and
narrative historians. But they will find the inscriptions more rewarding for what H.
Stuart Hughes has called ``retrospective cultural anthropology,' an analysis of the
cultural condition of a group of people at a certain place and time. The historian of \pinyinterm{shang}
China must constantly attempt to understand the sources in situ, as parts of an entire
cultural system with its own necessary rationale. The questions must be not only What
does the inscription say? but What function did it serve to say it? What did the record
keeper think he was doing?
Divination in its own right was a significant aspect of \pinyinterm{shang} life. To understand the
\pinyinterm{shang}, we must understand both the logic of their divination system and its social function.
It is anachronistic, and patronizing, to dismiss \pinyinterm{shang} pyromancy as mere superstition,
evidence that the \pinyinterm{shang} did not dare make up their minds for themselves.\footnote[89]{This view was common among Chinese scholars; see, e.g., \pinyinterm{hu-houxuan} (\citeyear{Hu1944aYintai}), p. 2a; \pinyinterm{dong-short} (\citeyear{Tung1965Chia}), p. 117.} %89
 To lament
that most of the divination charges are trivial is to misunderstand both the nature of
% source: scan 172, printed 154
\pinyinterm{shang} culture and the role of the historian.\footnote[90]{E.g., Soymié (\citeyear{Hu1959Yinputzu}), p. 278: the inscriptions ``notent les questions, des plus triviales aux plus graves (surtout les triviales) qui intéressaient les rois Yin.''} It is for the historian to discover what it
was that made these divinations nontrivial to the \pinyinterm{shang}, who expended so much effort
in preparing, cracking, and engraving the scapulas and plastrons; their ``triviality'' is a
clue of great significance if we would cross the cultural gap that separates us from Bronze
Age China. Cicero's comments, quoted as the epigraph to this book, should remind us
of the major role that divining-or forecasting-plays in our own culture. We all plan
for, and seek reassurance about, the future in ways that differ but express the central
concerns of our lives. Evans-Pritchard, to cite but one example of the naturalness of
divination, found that the poison oracle was ``as satisfactory a way of running my home
and affairs as any other I know of. Among Azande it is the only satisfactory way of life
because it is the only way of life they understand, and it furnishes the only arguments by
which they are wholly convinced and silenced.''\footnote[91]{Evans-Pritchard (\citeyear{EvansPritchard1937Witchcraft}), pp. 270, 313.} As he concluded, ``Magic and oracles
are more difficult for us to understand than most other primitive practices and are there-
fore more profitable a subject for study than customs which invite easy explanations.''

It is the primary assumption of this book, and of the studies that will follow, that the
divination inscriptions may be used to delineate accurately, though by no means com-
pletely, the ethos and world view of the \pinyinterm{shang} elite. The \pinyinterm{shang}, like the Azande, believed
in the explanatory, rational nature of divination; they knew the world worked this way.
\pinyinterm{shang} pyromancy is thus one of the clues which help us to define the constraining struc-
tures of the \pinyinterm{shang} mind, to characterize and explicate the ``culturally standardized
unreason''\footnote[92]{See Lévi-Strauss (\citeyear{LeviStrauss1969Raw}), p. 10; Kluckhohn and Leighton (\citeyear{Kluckhohn1970The}), p. 348.} which lay at the heart of \pinyinterm{shang} culture as it does all others. The great
questions of \pinyinterm{shang} history are not just political or economic: How were taxes collected?
What were the sources of \pinyinterm{shang} power and prosperity? They are also cultural: Who did
the \pinyinterm{shang} think they were? How did they conceive of themselves and their world? What
theological and metaphysical assumptions made sense of that world and their role in it?
and What significance did these conceptions have for the later development of Chinese
values? In the words of Marc Bloch, ``in the last analysis it is human consciousness which
is the subject-matter of history.''\footnote[93]{Bloch (\citeyear{Bloch1953Craft}), p. 151.} In answering for the \pinyinterm{shang} the human question, Who
am I? or, more accurately in \pinyinterm{shang} terms, perhaps, Who were, are, and will we be?---
scholars may, in small part, help us answer the question for ourselves. These are the
questions which touch us directly and deeply, and they are the questions to which the
graphs carved on bone and shell some three thousand years ago can still provide responses.
To study \pinyinterm{shang} conceptions of life is to extend our own.\footnote[94]{Cf. Winch (\citeyear{Winch1964Understanding}), esp. pp. 317-322.}
% source: scan 173, printed 155
In all this, the role of historiography is fundamental. Painstaking work is needed if
we are to see things clearly at a distance of three millennia. Much about the oracle-bone
inscriptions is not known or understood, much that we now know will change. I hope,
however, that the concerns introduced in this book will assist historians in their attempt
to see the \pinyinterm{shang} as they really were. For it was from their beginnings in the last half of the
second millennium B.C. that one of the greatest and most enduring of human cultures
arose. The origins of that culture, as we are humans too, concern us all.

\section[Methods of Transcription]{Methods of Transcription}
\label{sec:chapters-ch05:5.6}

% source: scan 165, printed 147
The translation by transcription of oracle bone graphs into modern Chinese graphs introduces (as already suggested (sec.~\ref{sec:chapters-ch03:3.5}% was: (sec. 3.5)
)) the risk of distortion. It is a classic case of traduttore, traditore, all the more insidious because likely to go unremarked.\footnote[59]{On some of the problems involved, see Serruys (\citeyear{Serruys1974The}), pp. 15-19, and, for a particularly instructive case of erroneous transcription, p. 101, n. 23.} %59
 Ideally, every inscription quoted should be presented as a rubbing and be accompanied by a tracing of the \pinyinterm{shang} characters, followed by a transcription into modern characters and then by a translation into a modern language. Such a method, in a work which cites several hundred inscriptions, would be impossibly cumbersome. Some scholars have chosen to write all the inscriptions they discuss in \pinyinterm{shang} form alone.\footnote[60]{E.g., \pinyinterm{shima-short} (\citeyear{Shima1958Inkyo}); Zhang Zongdong (\citeyear{Zhang1970Der}); Takashima (\citeyear{Takashima1973Negatives}).} %60
 I have followed their purist example by providing in the margin \pinyinterm{shima-short}'s oracle-bone transcriptions, taken from \booktitle{Sōrui} (sec.~\ref{sec:chapters-ch03:3.3.2}% was: (sec. 3.3.2)
), for some of the inscriptions translated in this book. For the scholar, steeped in \pinyintext{jiaguwen}, such transcriptions will generally suffice, at least as transcriptions (see the next paragraph).

Certain caveats should be applied to all forms of reproduction: one cannot always tell if the inscription is on bone or shell, one cannot tell its thickness (an important consideration in rejoining fragments; sec.~\ref{sec:chapters-ch05:5.7}% was: (sec. 5.7)
), one cannot tell what hollow forms or even graphs may have been on the back, one cannot even be sure that the shape of the entire fragment has been reproduced.\footnote[57]{The top of \pinyinterm{yicun} 888, for example, was cut off; \pinyinterm{hu-houxuan} (\citeyear{Hu1964Chiakuwen}), p. 138 and fig. 3, restores it. Other rubbings were cut in two; e.g., \pinyinterm{fuyin}, ``Ti'' 207 and ``Wen'' 40 (see the complete fragment at \pinyinterm{xubian} 1.45.3 or \pinyinterm{yicun} 914 [胡, op. cit., p. 156]); \pinyinterm{qianbian} 1.14.7 and 1.16.3 (rejoined at \pinyinterm{chaocun} 11). Even the \pinyinterm{yibian} rubbings do not always reproduce the whole fragment; e.g., the upper left quadrant of fragment F of \pinyinterm{pinbian} 395 is not reproduced in \pinyinterm{yibian} 3063. In his discussion of the oblong tablets into which oracle bones were occasionally sawn (sec.~\ref{sec:chapters-ch01:1.3.1}% was: (sec. 1.3.1)
), \pinyinterm{liu-yuanlin} (\citeyear{Liu1969Kuchien}), pp. 237-243, notes that the existence of such tablets had long been unrecognized because early books of rubbings did not always reproduce the full shape of the bone or shell fragment.} %57
 It goes without saying that scholars should work with the actual bones or shells when possible.\footnote[58]{It is probable that my discussion of \pinyinterm{pinbian} 12-21 (sec.~\ref{sec:chapters-ch03:3.7}% was: (sec. 3.7)
) would require revisions if I had the opportunity to examine the plastrons themselves. For examples of the dangers involved when the actual oracle bone is not studied, see \pinyinterm{yan-yiping} (\citeyear{Yen1961Chiaku}), pp. 516-520.} %58


\section[Reconstruction of Inscriptions]{Reconstruction of Inscriptions}
\label{sec:chapters-ch05:5.7}

% source: scan 167, printed 149
Many of the oracle bones excavated at \pinyinterm{xiaotun} were already in fragments when they came out of the ground; others were broken later by accident or to be sold separately. Rejoining these fragments, and thus the inscriptions, is a work of vital importance.\footnote[64]{On the method of rejoining fragments, see \pinyinterm{dong-short} (\citeyear{Tung1945Yin}), pt. 2, ch. 5, pp. 22a ff.; ch. 8, pp. 12a-13a; (\citeyear{Tung1954aYinli}), p. 120; (\citeyear{Tung1958Chiaku}), pp. 915-918; \pinyinterm{zhang-bingquan} (\citeyear{Zhang1956aPutzu}), pp. 242-243; Serruys (\citeyear{Serruys1974The}), pp. 15-19. Woodhead (\citeyear{Woodhead1967Study}), pp. 67-76, is also useful for the \pinyinterm{shang} scholar. Note that a small part of the rejoiner's work may involve rejoining fragments which were once whole, or rubbings which were subsequently cut in pieces. For example, \booktitle{Tieyun} 108.1, originally published as one fragment, was subsequently broken in two and published as \pinyinterm{xucun} 1.358 and \booktitle{Shitao} 2.92 (see \pinyinterm{zhuixin} 285). For rubbings, see n. 57.} %64
 It permits us to reconstitute whole divination sentences, to place them in the context of adjacent inscriptions, and even, on occasion, in the context of a series of inscriptions. Rejoinings may also serve, as already noted (sec.~\ref{sec:chapters-ch05:5.4.1}% was: (sec. 5.4.1)
, par. 1), to authenticate an inscription that was not scientifically excavated.\footnote[65]{See n. 45. The fact that \pinyinterm{jingjin} 1266 may be joined to \pinyinterm{pinbian} 24 (see \pinyinterm{pinbian} 24, kǎoshì, pp. 47-51) suggests either that the excavators of YH127 were careless and the piece was stolen or that the plastron was broken before the \pinyinterm{shang} put it in the ground. Cf. \pinyinterm{yan-yiping} (\citeyear{Yen1973Kuanyu}), pp. 2b, 6a, 6b-7a, for this and two similar cases.} %65
 Certain collections are composed entirely of rejoined fragments; and many other joinings have been made on an individual basis.\footnote[66]{The major published collections of rejoined pieces are \pinyinterm{chaocun}, \pinyinterm{zhuihe}, \pinyinterm{hebian}, and \pinyinterm{pinbian}; \pinyinterm{jiatu} (plates 1-58 at the rear of \pinyinterm{jiabian} kǎoshì) contains 221 rejoinings of \pinyinterm{jiabian} fragments. \pinyinterm{zhuixin} reproduces by rubbing and drawing many of these rejoinings (with the exception of those involving \pinyinterm{yibian} fragments) and many new ones; the preparation of an index would render this useful work more valuable. For a review of \pinyinterm{zhuixin} see Mickel (\citeyear{Mickel1976bThree}). Mickel (\citeyear{Mickel1974Reduplicated}), (\citeyear{Mickel1974aIndex}), (\citeyear{Mickel1976aIndex}) and (\citeyear{Mickel1977aIndex}) provide indexes to the rejoined (and reduplicated) fragments in \pinyinterm{jianshou}, \pinyinterm{qianbian}, \pinyinterm{jinghua}, \pinyinterm{fujia-title}, \pinyinterm{fuyin}, \pinyinterm{houbian}, \pinyinterm{xubian}, Kikkō, \pinyinterm{buci}, \booktitle{Tieyi}, \booktitle{Tieyu}, \booktitle{Tieyun}, \pinyinterm{cunzhen}, \pinyinterm{dong-short} (\citeyear{Tung1929aHsin}), \pinyinterm{yicun}, and 殷墟 (see table 9).} %66
 The work of rejoining, which may be likened to an attempt to put together hundreds of separate, yet deceptively similar, jigsaw puzzles from thousands of similar pieces, many of which are missing, some of which have lost their original shape (and some of which may even be rejoined in more than one way), may eventually be speeded up by the application of computer technology.\footnote[67]{For the correct rearrangement of the fragments wrongly rejoined in \pinyinterm{jiatu} 5, see n. 76.} %67


\pinyinterm{dong-zuobin}'s tongue-in-cheek claim that if one knows the kan-chih date, inscription type, and structure of the bone and shell, reconstruction is ``easier than putting a broken mirror together''\footnote[68]{\pinyinterm{dong-short} (\citeyear{Tung1954aYinli}), p. 120.} should remind us that any rejoining, especially if made only on the basis of rubbings, must be treated with caution. The rejoiner needs to be aware of the following pitfalls: rubbings of individual fragments may stretch or shrink to different degrees so that one rubbing will not match another, even though the original fragments do; it may be difficult to fit rubbings together since the curvature of the reconstituted bone or shell will differ from that of the rubbings made from individual fragments; differing pressures and conditions underground may have warped, broken, or discolored fragments from the same piece in different ways so that even the fragments themselves will not fit together except under pressure.\footnote[69]{\pinyinterm{zhou-hongxiang} (\citeyear{Chou1973Computer}).} %69


% source: scan 174, printed 156
Identification of the Inscribed
Turtle Shells of \pinyinterm{shang}
by James F. Berry
I.
Introduction
There have been a number of limited attempts at identifying the turtle remains from the
archaeological site at \pinyinterm{anyang}. Bing Zhi (\citeyear{Ping1930Notes}) described an extinct terrestrial tortoise, Testudo
anyangensis, based on an uninscribed shell from \pinyinterm{anyang}. Lindholm (\citeyear{Lindholm1931Uber}) concluded that the
turtle examined by Bing Zhi was, in fact, an aquatic turtle closely related to the living turtle Ocadia
sinensis, and he renamed the turtle Pseudocadia anyangensis. Pope (\citeyear{Pope1935Reptiles}), in his monograph on the
reptiles of China, accepted Lindholm's redescription and included a lengthy discussion of the
presumably extinct Pseudocadia.
In 1937 more of the \pinyinterm{anyang} turtle remains, including the specimen described as Pseudocadia
anyangensis, were examined by Bien. He concluded that the differences between Pseudocadia and
the living form Ocadia sinensis were not sufficient to justify the recognition of two species, and he
placed the name Pseudocadia anyangensis in synonymy with Ocadia sinensis. In addition, Bien was
first to recognize the remains of another living species, Chinemys (= Geoclemmys) reevesi, at the \pinyinterm{anyang} site. H. W. Wu (\citeyear{Wu1943Notes}) identified the largest plastron (\pinyinterm{pinbian} 184) as a third species, the
tortoise Testudo emys.
Several American authors have made reference to the turtles from the \pinyinterm{anyang} site. Pope
(cited in Carr (\citeyear{Carr1952Handbook})) referred to the presence of Pseudocadia at \pinyinterm{anyang}. Possibly unfamiliar with
the work of Lindholm and Bien, Auffenberg (\citeyear{Auffenberg1962Status}) commented on Bing Zhi's original description of
Testudo anyangensis, stating that the shell was definitely that of a freshwater turtle rather than that
of a terrestrial tortoise as Bing Zhi had suggested (the same conclusion reached earlier by Lindholm).
Auffenberg suggested that the \pinyinterm{anyang} turtle should be renamed Clemmys anyangensis (Bing Zhi).
McDowell (\citeyear{McDowell1964Partition}), in a taxonomic revision of the turtle family Testudinidae (the family Emydidae
of many authors), suggested that Testudo anyangensis Bing Zhi should be placed in synonymy with the
living species Mauremys mutica based on the greater likelihood of its having occurred in the area
near \pinyinterm{anyang}, and on the inadequacy of several characteristics used by Bing Zhi, Lindholm, and
Bien in distinguishing Mauremys from Ocadia. Mauremys mutica would thus represent the fourth
species represented at \pinyinterm{anyang}.
Ting Su (\citeyear{Su1969Shuo}) has made the only attempt to date to identify the \pinyinterm{anyang} turtle shells quantitatively. By comparing ratios of various seams between epidermal scutes on a limited sample of
